\documentclass[11pt]{article}

\usepackage[letterpaper,margin=1in,headheight=14pt]{geometry}
\usepackage[T1]{fontenc}
\usepackage[utf8]{inputenc}
\usepackage{lmodern}
\usepackage{microtype}
\usepackage{amsmath,amssymb,amsthm,mathtools,mathrsfs}
\usepackage{booktabs,longtable,tabularx,array}
\usepackage{xcolor}
\usepackage{enumitem}
\usepackage{fancyhdr}
\usepackage{titlesec}
\usepackage[hidelinks]{hyperref}
\usepackage[nameinlink,noabbrev]{cleveref}

\definecolor{ink}{HTML}{1F2937}
\definecolor{accent}{HTML}{1F5A7A}
\definecolor{soft}{HTML}{F3F6F8}
\hypersetup{
  pdftitle={Integrated Structures and Boundary Spectra in Sparse Prime Codings},
  pdfauthor={KEIO},
  colorlinks=true,
  linkcolor=accent,
  citecolor=accent,
  urlcolor=accent
}

\titleformat{\section}{\Large\bfseries\color{ink}}{\thesection}{0.7em}{}
\titleformat{\subsection}{\large\bfseries\color{ink}}{\thesubsection}{0.6em}{}
\titleformat{\subsubsection}{\normalsize\bfseries\color{ink}}{\thesubsubsection}{0.5em}{}
\setlist{itemsep=0.25em,topsep=0.5em}
\setlength{\parindent}{1.15em}
\setlength{\parskip}{0.15em}
\allowdisplaybreaks

\pagestyle{fancy}
\fancyhf{}
\fancyhead[L]{\small Integrated Structures and Boundary Spectra in Sparse Prime Codings}
\fancyhead[R]{\small KEIO}
\fancyfoot[C]{\thepage}

\newtheorem{theorem}{Theorem}[section]
\newtheorem{proposition}[theorem]{Proposition}
\newtheorem{lemma}[theorem]{Lemma}
\newtheorem{corollary}[theorem]{Corollary}
\newtheorem{claim}[theorem]{Claim}
\theoremstyle{definition}
\newtheorem{definition}[theorem]{Definition}
\newtheorem{conjecture}[theorem]{Conjecture}
\theoremstyle{remark}
\newtheorem{remark}[theorem]{Remark}

\newcommand{\PP}{\mathbb P}
\newcommand{\Ftwo}{\mathbb F_{2}}
\newcommand{\one}{\mathbf 1}
\newcommand{\eps}{\varepsilon}
\newcommand{\Palt}{\mathcal P_{\mathrm{alt}}}
\newcommand{\cP}{\mathfrak P}
\newcommand{\Dabs}{\mathcal D}
\newcommand{\status}[1]{\textnormal{\textcolor{accent}{\footnotesize[#1]}}}
\newcommand{\dd}{\,\mathrm d}
\DeclareMathOperator{\ord}{ord}
\DeclareMathOperator{\sgn}{sgn}
\DeclareMathOperator{\lcm}{lcm}
\DeclareMathOperator{\supp}{supp}
\DeclareMathOperator{\Tr}{Tr}
\DeclareMathOperator{\FPTr}{FPTr}
\DeclareMathOperator{\oddpart}{oddpart}

\begin{document}
\hypersetup{pageanchor=false}

\begin{titlepage}
\centering
\vspace*{1.15in}
{\huge\bfseries\color{ink} Integrated Structures and Boundary Spectra\par}
\vspace{0.12in}
{\huge\bfseries\color{ink} in Sparse Prime Codings\par}
\vspace{0.35in}
{\Large Quantitative Orthogonality, Exact Correlation Inversion,\par
Prime-Gap Spectral Models, and Gumbel Boundary Data\par}
\vspace{0.8in}
{\large KEIO\par}
\vspace{0.25in}
{\large Audited integrated manuscript, Version 3.0 (English)\par}
\vspace{0.1in}
{\large 2 September 2026\par}
\vfill
\begin{minipage}{0.86\textwidth}
\small
\textbf{Scope and provenance.}
This is the post-Version-2.0 audited revision of the English reconstruction of the
33-page Korean manuscript \emph{Integrated Structures in Sparse Prime Codings}.
It preserves the earlier results, corrects the anchored wheel-state congruence in the
Gilbreath section, and incorporates only later deductions for which a complete argument
or a precise literature input is supplied.  The main additions are a quantitative
prime-block transfer theorem for M\"obius and Liouville sums, exact all-lag inversion of
the signed prime-pair correlation, a critical rank-one model for the alternating
prime-index spectrum, and the Mellin natural boundary of the dyadic numerator series,
whose inverse-logarithmic coefficients are the moments of the standard Gumbel law.
Exact reformulations of open problems are explicitly separated from progress on them.
No claim of historical priority is made without a dedicated bibliographic search.
\end{minipage}
\vfill
\end{titlepage}

\hypersetup{pageanchor=true}
\pagenumbering{roman}

\begin{abstract}
Let $p_1=2<p_2<\cdots$ be the primes, $g_k=p_{k+1}-p_k$, and
$b_n=\pi(n)\bmod 2$.  The adjacent XOR identity
\[
  \one_{\PP}(n)=b_{n-1}\oplus b_n
\]
transports the cumulative prime-parity word, without loss, among symbolic dynamics,
signed correlations, doubling orbits, prime-gap continued fractions, and dyadic
numerator recurrences.  We make the factor-complexity bounds explicit:
\[
 \exp\!\left(\left(\frac1{16}-o(1)\right)(\log m)^2\right)
 \le p_b(m)\le \exp\!\left(C\frac{m}{\log m}\right),
\]
where the lower constant is optimal within the present Maynard--CRT shadow-counting
method.  The same comparison and lower bound hold for every modular prime clock
$\pi(n)\bmod d$.  Consequently these words have zero topological entropy but are not
morphic; their complex generating functions have the unit circle as a natural boundary,
and the associated formal series over $\Ftwo$ is transcendental.

Every bounded sequence of vanishing Ces\`aro total variation is orthogonal, in Ces\`aro
mean, to both the M\"obius and Liouville functions.  For arbitrary bounded prime-block
labels this is sharpened to the uniform rate
\[
 \sum_{n\le x}\mu(n)a_{\pi(n)},\quad
 \sum_{n\le x}\lambda(n)a_{\pi(n)}
 \ll \lVert a\rVert_\infty x\frac{\log_3x}{\log_2x},
\]
as well as an absolute total-variation estimate for the Mertens and Liouville sums over
individual prime blocks.  Fixed-lag correlations yield a transition hierarchy and a
prime-rooted two-wall local limit.  The finite diffraction measure recovers $\pi(N)$
exactly through a chordal Wasserstein distance; after the Bragg peak is removed, the
renormalized high-pass diffraction converges to Haar measure.  For the signed prime-pair
correlation $C_X(h)$ a discrete Green inversion recovers every overlap count and, in
particular, the global parity sum at even prime endpoints.  Phase tomography separately
recovers all counts classified by prime-index distance.  These are exact decoding
identities, not solutions of the twin-prime or global parity problems.

The gap $g_k$ is also decoded by the depth of a doubling orbit, the first partial
quotient of a Gauss orbit, and a continuant denominator ratio.  We prove
$\log q_n=n\log\log n+O(n)$, irrationality exponent $2$ for three prime-gap continued
fractions, and the Hausdorff-dimension bound $2/3$ for the closure of the Gauss tails.
The fixed-depth Gilbreath obstruction language is described by affine and transfer
systems.  A counterexample to the Version~2.0 fixed-start wheel state is given and the
state is corrected by recording the actual endpoint residues; the resulting anchored
right congruence is exact.

For the dyadic numerators
\[
 N_n=\sum_{j=1}^n(-1)^{j+1}2^{p_n-p_j},
\]
one integer $N_n$ autonomously recovers the entire preceding prime history, and
$(N_n)$ is not $P$-recursive.  The three lacunary support series $F_Q,F_H,F_N$ are
hypertranscendental; their nonzero algebraic specializations are transcendental and their
rational specializations are Liouville numbers.  The Dirichlet series
$Z_N(s)=\sum N_n^{-s}$ has the imaginary axis as a natural boundary, and as
$s\downarrow0$ its full inverse-logarithmic expansion has coefficients
$(-1)^m\Gamma^{(m)}(1)$, the moments of the standard Gumbel law.  Hence its boundary
data recovers $\gamma$ and all positive-integer zeta cumulants.

Finally, the alternating prime-index Dirichlet series $P_{\rm alt}$ is the relative trace
of a critical rank-one perturbation that changes the inverse odd-prime spectrum into the
inverse even-prime spectrum.  Its Jessen functions give the vague zero-measure limit on
$\Re s>1/2$; the support has exact right edge
\[
 \sigma_{\mathrm e}=1.779544653546994\ldots,
 \qquad \cP(\sigma_{\mathrm e})=2^{1-\sigma_{\mathrm e}},
\]
and that boundary line is zero-free.  We close by isolating what remains open: in
particular $\sum_{n\le x}(-1)^{\pi(n)}=o(x)$ and Erd\H{o}s's alternating prime series
are not proved here.
\end{abstract}

\medskip
\noindent\textbf{Keywords.}
Prime parity; factor complexity; M\"obius disjointness; Liouville function; natural
boundary; signed correlation; diffraction; rank-one perturbation; Gumbel moments;
bounded prime gaps; continued fractions; Gilbreath conjecture; automaticity;
almost-periodic Dirichlet series.

\tableofcontents
\newpage
\pagenumbering{arabic}

\section{Introduction and theorem map}

Sparse arithmetic data often admit many encodings: an infinite word, a real number, a
power series, an orbit, a continued fraction, or a difference table.  The main difficulty
is not constructing these encodings but determining which information survives the
translation and where genuinely new input is required.  This paper performs that audit
for the cumulative parity of the prime-counting function.

Write
\[
 p_1=2<p_2<\cdots,\qquad g_k=p_{k+1}-p_k,
\]
and define
\begin{equation}\label{eq:basic-defs}
 \chi_n:=\one_{\PP}(n),\qquad b_0:=0,\qquad
 b_n:=\pi(n)\pmod 2,\qquad \sigma_n:=(-1)^{b_n}.
\end{equation}
The organizing identity is
\begin{equation}\label{eq:xor}
 \chi_n=b_{n-1}\oplus b_n.
\end{equation}
It is elementary, but it makes the translation dictionary exact.  The word $b$ and its
coefficient sequence are recorded in OEIS; our concern is the network of consequences,
not priority for the underlying digits.

The main additions in this post-Version-2.0 revision are as follows.
\begin{enumerate}[label=\textup{(\roman*)}]
 \item The Maynard--CRT complexity amplification is made explicit with constant
 $1/16$, optimal for the present shadow-counting argument, and is extended to the
 modular prime clocks $\pi(n)\bmod d$.
 \item The qualitative sparse-transition principle is complemented by a uniform
 $x\log_3x/\log_2x$ M\"obius--Liouville bound for arbitrary prime-block labels and by
 an absolute block-variation theorem.
 \item Prime-rooted wall limits, reciprocal occurrence laws, exact diffraction recovery,
 a renormalized Haar limit, all-lag Green inversion, and phase tomography refine the
 correlation dictionary without asserting prime-pair infinitude or global parity
 cancellation.
 \item The Gauss-tail closure is proved to have Hausdorff dimension at most $2/3$.
 The fixed-start wheel state in the Gilbreath language is corrected by adjoining actual
 endpoint residues, and a counterexample explains why the earlier relative mask was
 insufficient.
 \item A single dyadic numerator $N_n$ is shown to recover the full preceding prime
 history autonomously; OEIS A122150 is not $P$-recursive.  The three support series are
 strengthened to hypertranscendence, with arithmetic Picard fibres for $F_N$.
 \item The new Dirichlet series $Z_N(s)=\sum N_n^{-s}$ has natural boundary
 $\Re s=0$ and a complete Gumbel-moment boundary expansion.  The Liouville values
 $F_N(r)$ also yield algebraic-independence statements with $\pi$, the even zeta values,
 and infinitely many of those boundary moments.
 \item The odd and even inverse-prime spectra are connected by a critical rank-one
 perturbation whose relative trace is $\Palt$.  Generalized Jessen theory then gives the
 limiting zero measure throughout $\Re s>1/2$.
 \item The global parity frontier is sharpened by exact correlation reconstruction,
 large-gap excursions, and convergence at every nonzero Fourier boundary frequency;
 the unweighted zero-frequency problem remains explicitly open.
\end{enumerate}

Five corrections to the earlier version are built into the notation and statements.  We
write $2C_{PP}$ literally rather than introducing a tilde; the Gauss-tail indexing is
$G^k(\alpha_P)$, not $G^{k+1}(\alpha_P)$; the first exponent of each of $F_H$ and $F_N$
has multiplicity two; and previously suggested positive-density claims for
$\pi(n)\bmod2$ are replaced by the stronger published results of Chung--Li and Alboiu.
In addition, the fixed-start Gilbreath quotient now records actual endpoint residues;
the former relative mask survives only for the distinct existential-translate problem.

We use the following provenance tags only as reading aids:
\begin{center}
\begin{tabular}{@{}cl@{}}
\toprule
\status{E} & exact or elementary argument given here;\\
\status{A} & standard analytic theorem, with the input identified;\\
\status{L} & substantial theorem imported from the literature;\\
\status{C} & conditional statement, with its hypothesis displayed.\\
\bottomrule
\end{tabular}
\end{center}
The tags describe proof dependence, not novelty or journal value.

\section{Prime-parity coding and the doubled constant}

\subsection{Two finite differences}

Set $a_n=b_n-b_{n-1}$ as an integer difference.  Then
\begin{equation}\label{eq:signed-prime-difference}
 a_n=
 \begin{cases}
   (-1)^{k+1},&n=p_k,\\
   0,&n\notin\PP.
 \end{cases}
\end{equation}
Consequently, define
\[
 B(z):=\sum_{n\ge1}b_nz^n,
 \qquad A(z):=\sum_{k\ge1}(-1)^{k+1}z^{p_k},
\]
and let $\overline B$ denote the coefficientwise reduction of $B$ modulo $2$.  We have
the exact identities
\begin{equation}\label{eq:GF-dictionary}
 A(z)=(1-z)B(z),
 \qquad
 (1+z)\overline B(z)=\sum_{p\in\PP}z^p
 \quad\text{in }\Ftwo[[z]].
\end{equation}
The first identity controls complex-analytic and correlation questions; the second
transports automaticity and formal algebraicity.

\subsection{The constant, its digits, and the doubling detector}

Define the parity prime constant
\begin{equation}\label{eq:CPP-def}
 C_{PP}:=\sum_{k\ge1}(-1)^{k+1}2^{-p_k}.
\end{equation}
Evaluating the first identity in \eqref{eq:GF-dictionary} at $z=1/2$ gives
\begin{equation}\label{eq:2CPP-binary}
 2C_{PP}=\sum_{n\ge1}b_n2^{-n}=0.b_1b_2b_3\cdots{}_2.
\end{equation}
Thus $C_{PP}=0.0b_1b_2\cdots{}_2$.  This one-place shift matters when comparing the
constant with OEIS entries A071986, A122150 and A122153: A071986 supplies the binary
digits of $2C_{PP}$, rather than of $C_{PP}$ under the standard radix-point convention.

\begin{corollary}[Irrationality and non-normality]\label{cor:CPP-irrational}
\status{E}
Both $C_{PP}$ and $2C_{PP}$ are irrational and non-normal in base $2$.
\end{corollary}

\begin{proof}
If $b$ were eventually periodic, then its adjacent XOR word would be eventually periodic.
The latter is impossible because the primes are infinite and have arbitrarily long gaps.
Hence the binary expansion in \eqref{eq:2CPP-binary} is not eventually periodic, proving
irrationality.  The complexity upper bound in \cref{thm:complexity-upper} below is
$p_b(m)=\exp(O(m/\log m))=o(2^m)$, so for all large $m$ at least one binary block of
length $m$ never occurs.  A base-$2$ normal number contains every such block.  Adding or
removing the initial zero digit does not change this conclusion for $C_{PP}$.
\end{proof}

Let $T(x)=2x\pmod1$ and $R=[1/4,3/4)$.  Reading the first two digits after $n$ shifts
shows that, for every $n\ge0$,
\begin{equation}\label{eq:doubling-detector}
 \one_R\!\left(T^n(2C_{PP})\right)=\chi_{n+2}.
\end{equation}
The visit times to $R$ are therefore $p_k-2$, and successive return times are precisely
the gaps $g_k$.  This is an exact coordinate change, not a new primality test.

\section{Factor complexity, morphisms, automata, and natural boundaries}

Let $p_{\PP}(m)$ and $p_b(m)$ denote the number of distinct length-$m$ factors of the
prime characteristic word and of $b=b_0b_1\cdots$, respectively.

\begin{theorem}[Exact comparison]\label{thm:complexity-comparison}
\status{E}
For every $m\ge1$,
\begin{equation}\label{eq:complexity-comparison}
 p_{\PP}(m)\le p_b(m+1)\le2p_{\PP}(m).
\end{equation}
\end{theorem}

\begin{proof}
Adjacent XOR sends every factor of $b$ of length $m+1$ to a unique prime-indicator
factor of length $m$.  Conversely, an indicator factor and the initial bit of a lift
determine all remaining bits by cumulative XOR.  There are at most two choices for that
initial bit.
\end{proof}

\subsection{A sieve upper bound}

We use the following standard uniform form of the one-dimensional Selberg upper-bound
sieve; see, for example, \cite[Chapter 6]{IK}.

\begin{lemma}[Uniform omitted-residue sieve]\label{lem:uniform-sieve}
\status{A}
Put $z=m^{1/100}$.  For each prime $\ell\le z$ choose an arbitrary residue
$r_\ell\pmod\ell$, and set
\[
 \mathcal R(\mathbf r):=
 \{1\le j\le m:j\not\equiv r_\ell\pmod\ell\text{ for every }\ell\le z\}.
\]
Uniformly in the residue choices,
\[
 |\mathcal R(\mathbf r)|\ll \frac{m}{\log m}.
\]
\end{lemma}

\begin{proof}
The local density removed at $\ell$ is $1/\ell$, independently of the selected residue.
The dimension-one Selberg upper-bound sieve at level $z$ therefore gives
\[
 |\mathcal R(\mathbf r)|
 \ll m\prod_{\ell\le z}\left(1-\frac1\ell\right)
 \asymp \frac{m}{\log z}\asymp\frac{m}{\log m}.
\]
The Chinese remainder theorem makes the remainder estimates uniform in $\mathbf r$.
\end{proof}

\begin{theorem}[Subexponential upper bound]\label{thm:complexity-upper}
\status{A}
There is an absolute constant $C>0$ such that
\begin{equation}\label{eq:complexity-upper}
 p_{\PP}(m),p_b(m)\le
 \exp\!\left(C\frac{m}{\log m}\right).
\end{equation}
\end{theorem}

\begin{proof}
Consider a length-$m$ prime-indicator block starting at $N>z$.  For every
$\ell\le z$, its prime positions omit at least one residue class modulo $\ell$: the unique
position congruent to $-N\pmod\ell$ would otherwise represent a multiple of $\ell$ larger
than $\ell$.  Hence the set of $1$-positions lies in some $\mathcal R(\mathbf r)$.
There are
\[
 \prod_{\ell\le z}\ell=\exp(\vartheta(z))=\exp(O(z))
\]
possible residue vectors and, for each, at most
$2^{|\mathcal R(\mathbf r)|}$ subsets.  By \cref{lem:uniform-sieve} and
$z=o(m/\log m)$, the product is $\exp(O(m/\log m))$.  Starts $N\le z$ contribute only
$O(z)$ further blocks.  Transfer to $b$ by \cref{thm:complexity-comparison}.
\end{proof}

\subsection{Maynard--CRT amplification}

The upper bound has an unconditional superpolynomial counterpart.  The construction is
included because the CRT admissibility check is essential.

\begin{theorem}[Explicit complexity amplification]\label{thm:maynard-complexity}
\status{L+A}
As $L\to\infty$,
\begin{equation}\label{eq:maynard-lower}
 \log p_{\PP}(L),\ \log p_b(L)
 \ge\left(\frac1{16}-o(1)\right)(\log L)^2.
\end{equation}
The constant $1/16$ is optimal within the present Maynard--shadow-counting
argument.
\end{theorem}

\begin{proof}
Let $T$ be the set of primes in $(L/2,L]$ and write $K=|T|\asymp L/\log L$.
For large $L$ we have $K<L/2$.  The offset set $T$ is admissible: if $\ell\le K$ is
prime, the residue $0$ is absent because every member of $T$ exceeds $\ell$; if
$\ell>K$, then $K$ offsets cannot cover all $\ell$ residue classes.

Set $k=\lfloor\sqrt K\rfloor$.  For every $k$-element subset $U\subset T$, choose
distinct primes $q_a>L$ for all $a\in[1,L]\setminus U$.  The Chinese remainder theorem
gives $N_U$ with
\[
 N_U\equiv-a\pmod{q_a}\qquad(a\in[1,L]\setminus U),
\]
and put $M_U=\prod_{a\notin U}q_a$.  Consider the common-slope linear forms
\[
 \mathcal L_u(t):=M_Ut+N_U+u\qquad(u\in U).
\]
They remain admissible.  Indeed, if a prime $q$ divides $M_U$, then $q=q_a$ for one
$a\notin U$ and $N_U+u\equiv u-a\not\equiv0\pmod q$, because
$0<|u-a|<L<q$.  If $q\nmid M_U$, multiplication by $M_U^{-1}$ preserves the number of
occupied root classes, and admissibility follows from that of $U\subset T$.

Maynard's variational estimate
$M_k>\log k-2\log\log k-2$ and his criterion $M_k>4m$
\cite{MaynardSmallGaps} show that, for every $U$ and infinitely many $t$, at least
\[
 r=\left(\frac14-o(1)\right)\log k
\]
of the values $\mathcal L_u(t)$ are prime.  Taking $t$ large also ensures that every
$N_U+M_Ut+a$ with $a\notin U$ is a composite multiple of $q_a$.  Thus each $U$
produces a length-$L$ prime pattern $S\subset U$ with $|S|\ge r$.

A fixed pattern $S$ of size at least $r$ can arise from at most
$\binom{K-r}{k-r}$ choices of $U$.  Consequently the number of distinct patterns is at
least
\[
 \frac{\binom Kk}{\binom{K-r}{k-r}}
 =\prod_{j=0}^{r-1}\frac{K-j}{k-j}
 \ge \exp\!\left(\left(\frac1{16}-o(1)\right)(\log K)^2\right)
 =\exp\!\left(\left(\frac1{16}-o(1)\right)(\log L)^2\right).
\]
The bound for $b$ follows from \eqref{eq:complexity-comparison}, after an immaterial
shift of $L$.

More generally, if $k=K^\alpha$ with $0<\alpha<1$, the same computation gives the
constant $\frac14\alpha(1-\alpha)$.  Its maximum is $1/16$, attained at
$\alpha=1/2$.
\end{proof}

\begin{corollary}[Modular prime clocks]\label{cor:modular-clocks}
\status{E+A}
For every integer $d\ge2$, let
\[
 c_n^{(d)}:=\pi(n)\pmod d
\]
and denote its factor complexity by $p_d(m)$.  Then
\begin{equation}\label{eq:modular-clock-comparison}
 p_{\PP}(m)\le p_d(m+1)\le d\,p_{\PP}(m),
\end{equation}
so
\[
 \log p_d(m)\ge
 \left(\frac1{16}-o(1)\right)(\log m)^2,
 \qquad
 \log p_d(m)=O(m/\log m).
\]
Every invariant measure on the shift orbit closure of $c^{(d)}$ is supported on
constant fixed points.
\end{corollary}

\begin{proof}
Taking adjacent differences modulo $d$ recovers the prime indicator, and an indicator
factor together with its first clock value has at most $d$ lifts.  This proves
\eqref{eq:modular-clock-comparison}.  The changes of $c^{(d)}$ occur only at primes,
whose upper Banach density is zero; the orbit-closure argument of
\cref{rem:upper-Banach} therefore forces every invariant measure onto constant points.
The assertion does not claim that every one of the $d$ constant points belongs to the
orbit closure.
\end{proof}

\begin{corollary}[Non-morphicity]\label{cor:nonmorphic}
\status{A}
Neither the prime characteristic word nor the cumulative parity word $b$ is morphic.
\end{corollary}

\begin{proof}
Every morphic word has factor complexity $O(n^2)$ by the general complexity theorem of
Devyatov \cite{Devyatov}.  The lower bound \eqref{eq:maynard-lower} exceeds every fixed
power of $L$.
\end{proof}

\begin{theorem}[Conditional optimal order]\label{thm:conditional-complexity}
\status{C}
Assume the prime $k$-tuple conjecture for every finite admissible tuple.  Then
\[
 p_{\PP}(L),p_b(L)=
 \exp\!\left(\Theta\!\left(\frac{L}{\log L}\right)\right).
\]
\end{theorem}

\begin{proof}
Use the same admissible $T$.  For each prescribed subset $S\subset T$, force every
position in $[1,L]\setminus S$ composite by distinct moduli as above.  The remaining
forms indexed by $S$ are admissible, so the conjecture makes all of them prime
simultaneously.  This realizes $2^K=\exp(\Omega(L/\log L))$ distinct patterns.  The
reverse inequality is \cref{thm:complexity-upper}.
\end{proof}

\subsection{Entropy, automaticity, and two natural boundaries}

\begin{corollary}\label{cor:zero-entropy}
\status{A}
The word $b$ is aperiodic and
\[
 m+1\le p_b(m)\le\exp(O(m/\log m)),
 \qquad h_{\mathrm{top}}(b)=0.
\]
\end{corollary}

\begin{proof}
Aperiodicity follows from \eqref{eq:xor}, the infinitude of primes, and arbitrarily long
prime gaps.  Morse--Hedlund gives $p_b(m)\ge m+1$; the upper bound gives zero entropy.
\end{proof}

\begin{theorem}[Automaticity and formal transcendence]\label{thm:automaticity}
\status{L+E}
For no integer base $q\ge2$ is $b$ a $q$-automatic word.  Moreover
\[
 \overline B(z)=\sum_{n\ge1}b_nz^n\in\Ftwo[[z]]
\]
is transcendental over $\Ftwo(z)$.
\end{theorem}

\begin{proof}
Automatic words are closed under shifts and coordinatewise finite-alphabet operations.
If $b$ were $q$-automatic, \eqref{eq:xor} would make the prime characteristic word
$q$-automatic, contradicting the nonrecognizability theorem for primes
\cite{HartmanisShank,AlloucheShallit}.  For $q=2$, Christol's theorem
\cite{ChristolEtAl} identifies algebraic power series over $\Ftwo(z)$ with $2$-automatic
coefficient sequences, proving the second assertion.
\end{proof}

\begin{theorem}[Natural boundaries]\label{thm:AB-natural-boundary}
\status{A}
The functions $A(z)$ and $B(z)$ have radius of convergence $1$ and the unit circle as a
natural boundary.  Hence both are transcendental over $\mathbb C(z)$ and are not
D-finite.
\end{theorem}

\begin{proof}
Their coefficients lie in finite sets of integers and have infinitely many nonzero terms,
so the radius is $1$.  A rational power series with coefficients in a finite set is
eventually periodic.  Rationality of $B$ would make $b$ eventually periodic; rationality
of $A$ would make its signed prime support eventually periodic.  Both are impossible.
The P\'olya--Carlson dichotomy \cite{Polya,Carlson} therefore gives the natural boundary.
An algebraic function, and more generally a D-finite function, has only finitely many
singularities in the finite plane, so neither class can contain $A$ or $B$.
\end{proof}

\begin{remark}
Formal or functional transcendence does not imply the transcendence of the particular
value $B(1/2)=2C_{PP}$.  The latter remains open.
\end{remark}

\section{Sparse transitions and M\"obius--Liouville orthogonality}

The prime-parity word has zero entropy for a quantitative combinatorial reason, but it
also has a much stronger measure-theoretic degeneracy: it changes only at a zero-density
set.  The following criterion extracts the multiplicative consequence.  The deep input is
the measure-complexity theorem of Huang, Wang and Ye \cite{HWY}; the deduction below is
short and should be read as an application of that theorem.  Wei's work on asymptotically
periodic functions \cite{Wei} is closely related; in particular it covers block-constant
sequences whose successive block lengths tend to infinity.  Our density criterion permits
recurring short clusters and therefore applies directly to prime boundaries.

For a bounded complex sequence $u=(u_n)_{n\ge1}$ put
\[
 V_u(N):=\sum_{n=1}^{N}|u_{n+1}-u_n|.
\]

\begin{theorem}[Vanishing Ces\`aro variation]\label{thm:variation-orthogonality}
\status{L+E}
If $u$ is bounded and $V_u(N)=o(N)$, then
\begin{equation}\label{eq:mu-lambda-orthogonality}
 \sum_{n\le N}\mu(n)u_n=o(N),
 \qquad
 \sum_{n\le N}\lambda(n)u_n=o(N).
\end{equation}
Here $\mu$ and $\lambda$ denote the M\"obius and Liouville functions.
\end{theorem}

\begin{proof}
Let $K\subset\mathbb C$ be the compact closure of the values of $u$.  Extend $u$ by a
constant negative tail to a point $x\in K^{\mathbb Z}$, let $S$ be the bilateral shift,
and take the orbit closure $X$ of $x$.  Given any sequence $N_j\to\infty$, pass to a
subsequence for which
\[
 \rho_j:=\frac1{N_j}\sum_{n=1}^{N_j}\delta_{S^nx}
 \quad\Longrightarrow\quad \rho
\]
weakly.  The limit is $S$-invariant.  For the continuous function
$\Psi(y)=|y_1-y_0|$,
\[
 \int_X\Psi\,\dd\rho
 =\lim_j\frac{V_u(N_j)}{N_j}=0.
\]
Thus $y_1=y_0$ for $\rho$-almost every $y$; invariance and a countable intersection
show that $\rho$ is supported on constant sequences.  The shift is therefore the identity
$\rho$-almost everywhere.  Such a measure has bounded, hence subpolynomial, measure
complexity.  Huang--Wang--Ye's pointwise criterion
\cite[Theorem~1.1$'$]{HWY}, applied to the coordinate observable $y\mapsto y_0$, gives
\[
 \frac1{N_j}\sum_{n\le N_j}\mu(n)u_n\longrightarrow0
\]
along this subsequence.  A bad-subsequence contradiction proves the full M\"obius limit.

For Liouville, use the exact square-divisor identity
\[
 \lambda(n)=\sum_{d^2\mid n}\mu(n/d^2).
\]
For each fixed $d$, the dilated sequence $u^{(d)}_m=u_{d^2m}$ also has vanishing
Ces\`aro variation, since
\[
 \sum_{m<H}|u_{d^2(m+1)}-u_{d^2m}|
 \le\sum_{r\le d^2H+O(d^2)}|u_{r+1}-u_r|=o(H).
\]
Hence the contribution of every fixed $d$ is $o_d(N/d^2)$.  For fixed $D$, the terms
$d\le D$ total $o(N)$, while the remaining terms are bounded by
\[
 \|u\|_\infty N\sum_{d>D}d^{-2}=O(\|u\|_\infty N/D).
\]
Let $N\to\infty$ and then $D\to\infty$.
\end{proof}

\begin{remark}[Orbit-closure form]\label{rem:upper-Banach}
If the stronger uniform variation condition
\[
 \lim_{H\to\infty}\sup_m\frac1H
 \sum_{m\le n<m+H}|u_{n+1}-u_n|=0
\]
holds, then every invariant measure on the entire orbit closure is supported on constant
fixed points.  Huang--Wang--Ye's system-wide theorem \cite[Theorem~1.2]{HWY} gives
M\"obius disjointness for the whole subshift.  For a finite-valued word it is enough that
its change set have upper Banach density zero.  Ordinary density zero suffices for the
pointwise \cref{thm:variation-orthogonality}, but not for this orbit-closure statement.
\end{remark}

The primes themselves have upper Banach density zero.  Indeed, for
$Q_z=\prod_{p\le z}p$, primes larger than $z$ occupy only the $\varphi(Q_z)$ reduced
classes modulo $Q_z$; uniformly over intervals of length $H$ their upper density is at
most $\varphi(Q_z)/Q_z+O_z(1/H)$, which tends to zero as first $H$ and then $z$ tend to
infinity.

\begin{corollary}[Arbitrary labels on prime-gap blocks]\label{cor:prime-block-labels}
\status{L+E}
For every bounded complex sequence $(a_k)_{k\ge0}$,
\begin{equation}\label{eq:any-prime-label}
 \sum_{n\le N}\mu(n)a_{\pi(n)}=o(N),\qquad
 \sum_{n\le N}\lambda(n)a_{\pi(n)}=o(N).
\end{equation}
The same is true with $\pi(n-1)$ in place of $\pi(n)$.
\end{corollary}

\begin{proof}
The sequence $a_{\pi(n-1)}$ changes only at primes.  Its total variation is at most
$2\sup_k|a_k|\,\pi(N)=o(N)$, so \cref{thm:variation-orthogonality} applies.  Replacing
$\pi(n-1)$ by $\pi(n)$ changes only the prime-indexed terms, whose total contribution is
$O(\pi(N))=o(N)$.
\end{proof}

Taking $a_k=(-1)^k$ yields, unconditionally,
\begin{equation}\label{eq:prime-parity-mu-lambda}
 \sum_{n\le N}\mu(n)(-1)^{\pi(n)}=o(N),
 \qquad
 \sum_{n\le N}\lambda(n)(-1)^{\pi(n)}=o(N).
\end{equation}
Every fixed finite-window coding
$\Phi(\sigma_{n+h_1},\ldots,\sigma_{n+h_r})$, and in particular every finite shift
product, changes only on a finite union of shifts of the primes.  It therefore satisfies
the same two orthogonality conclusions.

Two illustrative applications require no information about the underlying prime race.
If $\pi_2(n)$ counts twin-prime pairs with lower endpoint at most $n$, then
$(-1)^{\pi_2(n)}$ changes on a subset of the primes; both its $\mu$- and
$\lambda$-correlations are $o(N)$ whether the twin primes are finite or infinite.
Similarly,
\[
 \sgn\bigl(\pi(n;4,3)-\pi(n;4,1)\bigr)
\]
changes only at primes, so it too is orthogonal to both multiplicative functions.  These
are sparse-boundary consequences, not results on twin-prime infinitude or the limiting
Chebyshev bias.

\subsection{A quantitative prime-block transfer theorem}

The qualitative dynamical argument above is uniform in the labels but supplies no
rate.  Short-interval cancellation for multiplicative functions gives a complementary
quantitative statement.  We write $\log_j$ for the $j$-fold iterated logarithm.

\begin{theorem}[Quantitative prime-block transfer]
\label{thm:quant-prime-block}
\status{L+E}
Uniformly for every bounded complex sequence $(a_k)$,
\begin{equation}\label{eq:quant-prime-block}
 \sum_{n\le x}\mu(n)a_{\pi(n)},\quad
 \sum_{n\le x}\lambda(n)a_{\pi(n)}
 \ll \|a\|_\infty x\frac{\log_3x}{\log_2x}.
\end{equation}
If $f$ is either $\mu$ or $\lambda$ and $F(y)=\sum_{n\le y}f(n)$, then
\begin{equation}\label{eq:quant-block-variation}
 \sum_{p_{k+1}\le x}
 \left|F(p_{k+1}-1)-F(p_k-1)\right|
 \ll x\frac{\log_3x}{\log_2x}.
\end{equation}
The same bounds hold for $\pi(n-1)$, fixed finite-window codings, and fixed finite
shift products, with constants depending on the window.
\end{theorem}

\begin{proof}
It is enough to work on a dyadic interval $[X,2X]$.  Put
\[
 H=\left\lfloor\frac{\log X}{(\log_2X)^4}\right\rfloor,
 \qquad
 G=\left\lfloor\frac{\log X}{(\log_2X)^2}\right\rfloor,
 \qquad
 W_H(t)=\sum_{t<n\le t+H}f(n).
\]
The total length of prime-gap blocks of length at most $G$ is
$O(XG/\log X)$.  If a block $I$ has length greater than $G$, sliding an
$H$-window through its interior gives
\begin{equation}\label{eq:sliding-window-block}
 \left|\sum_{n\in I}f(n)\right|
 \le \frac1H
 \sum_{[t+1,t+H]\subset I}|W_H(t)|+O(H).
\end{equation}
There are $O(X/G)$ long blocks, so their total endpoint error is $O(XH/G)$.

The mean-square short-interval estimate of Matom\"aki--Radziwi\l\l--Tao and
Cauchy--Schwarz give
\[
 \frac1H\sum_{X\le t\le2X}|W_H(t)|
 \ll X\left(
 \sqrt{e^{-M(f;X)}M(f;X)}
 +\frac{\log\log H}{\log H}
 +\frac1{(\log X)^{1/100}}
 \right),
\]
where $M(f;X)$ is the pretentious-distance parameter in
\cite[Appendix~A, Theorem~A.1]{MRTChowla}.  Since
$\mu(p)=\lambda(p)=-1$, the bound
\cite[(1.12)]{MRTChowla} gives
\[
 M(\mu;X),M(\lambda;X)\ge\left(\frac13-o(1)\right)\log_2X.
\]
Thus the first and third terms are $O((\log X)^{-c})$ for some absolute $c>0$,
whereas $\log H\sim\log_2X$ and $\log\log H\sim\log_3X$.  Hence the preceding
display is $\ll X\log_3X/\log_2X$.  Combining this with the short-block contribution and
\eqref{eq:sliding-window-block} proves the same bound for the total absolute sum of
the multiplicative block sums.  Multiplication by arbitrary labels of modulus at most
$\|a\|_\infty$ proves \eqref{eq:quant-prime-block}; choosing on every block the
conjugate phase of its sum proves \eqref{eq:quant-block-variation}.  Dyadic summation
completes the proof.  A fixed window changes only $O(1)$ neighboring block boundaries,
so the final assertion follows by the same decomposition.
\end{proof}

\subsection{Weighted Mertens sums along zero-density sequences}

Let
\[
 M(x):=\sum_{n\le x}\mu(n),\qquad
 L(x):=\sum_{n\le x}\lambda(n).
\]

\begin{corollary}[Sparse sampling theorem]\label{cor:sparse-Mertens}
\status{L+E}
Let $1\le t_1<t_2<\cdots$ be infinite with
$\#\{k:t_k\le x\}=o(x)$.  If $(c_k)$ has uniformly bounded partial sums
$C_K=\sum_{k\le K}c_k$, then
\begin{equation}\label{eq:sparse-Mertens}
 \sum_{k\le K}c_kM(t_k)=o(t_K),
 \qquad
 \sum_{k\le K}c_kL(t_k)=o(t_K).
\end{equation}
\end{corollary}

\begin{proof}
Set $a_0=0$, $a_k=-C_k$, and
\[
 u_n=a_{A(n)},\qquad A(n):=\#\{j:t_j<n\}.
\]
The word $u$ is bounded and changes only at the zero-density set $\{t_k\}$.  At
$N_K=t_K+1$, discrete summation by parts gives the exact identities
\begin{align*}
 \sum_{n\le N_K}\mu(n)u_n
 &=M(N_K)a_K+\sum_{k\le K}c_kM(t_k),\\
 \sum_{n\le N_K}\lambda(n)u_n
 &=L(N_K)a_K+\sum_{k\le K}c_kL(t_k).
\end{align*}
The left sides are $o(N_K)$ by \cref{thm:variation-orthogonality}; also
$M(N)=o(N)$ by the prime number theorem, while $L(N)=o(N)$ follows from the Liouville
part of that theorem with $u\equiv1$.  Since $a_K$ is bounded and $K=o(t_K)$, the result
follows.
\end{proof}

In particular, for every zero-density increasing sequence,
\begin{equation}\label{eq:alternating-Mertens}
 \sum_{k\le K}(-1)^kM(t_k)=o(t_K),\qquad
 \sum_{k\le K}(-1)^kL(t_k)=o(t_K).
\end{equation}
This includes $t_k=p_k$.  More generally all mean-zero periodic weights and all
nonprincipal periodic characters are covered.

There is also an absolute block-cancellation consequence.  Choose on each interval
$(t_k,t_{k+1}]$ the phase conjugate to the corresponding M\"obius block sum and apply
\cref{thm:variation-orthogonality} to the resulting bounded block label.  Then
\begin{equation}\label{eq:block-absolute-cancellation}
 \sum_{k<K}|M(t_{k+1})-M(t_k)|=o(t_K),
\end{equation}
and analogously for $L$.  For prime blocks, \cref{thm:quant-prime-block} supplies the
explicit rate in \eqref{eq:quant-block-variation}; for a general zero-density sequence
the assertion remains qualitative.

\section{Fixed-lag correlations and transition strata}

For fixed $h\ge1$ define
\[
 D_h(N):=\#\{n\le N:b_{n+h}\ne b_n\},
 \qquad
 \rho_N(h):=\frac1N\sum_{n\le N}\sigma_n\sigma_{n+h}.
\]

\begin{theorem}[Fixed-lag asymptotic]\label{thm:fixed-lag}
\status{A}
For every fixed $h$,
\begin{align}
 D_h(N)&=h\pi(N)+O_h\!\left(\frac{N}{\log^2N}\right),\label{eq:Dh}\\
 \rho_N(h)&=1-\frac{2h}{\log N}+O_h(\log^{-2}N).\label{eq:rhoh}
\end{align}
\end{theorem}

\begin{proof}
Put $X_h(n)=\pi(n+h)-\pi(n)$.  The bits disagree precisely when $X_h(n)$ is odd, and
endpoint counting gives
\[
 \sum_{n\le N}X_h(n)=h\pi(N)+O_h(1).
\]
On a window containing at most one prime, $X_h(n)$ equals its oddness indicator.  A
discrepancy requires a pair of primes among $O_h(1)$ fixed offset pairs.  The Selberg
prime-pair upper bound makes the number of such windows
$O_h(N/\log^2N)$, proving \eqref{eq:Dh}.  Finally
$\sigma_n\sigma_{n+h}=1-2\one_{b_n\ne b_{n+h}}$ and the prime number theorem give
\eqref{eq:rhoh}.
\end{proof}

For every fixed integer $h$, $\rho_N(h)\to1$.  Hence the Herglotz measure associated
with these Besicovitch autocorrelation coefficients is $\delta_0$.  This does not assert
that the full orbit closure is uniquely ergodic; \cref{rem:upper-Banach} instead describes
all of its invariant measures, which are mixtures of constant fixed points.

For a binary word $w=w_1\cdots w_m$ define its transition count
\[
 \tau(w):=\#\{1\le j<m:w_j\ne w_{j+1}\},
\]
and let $A_w(N)$ count its occurrences starting at positions at most $N$.

\begin{theorem}[Transition-count hierarchy]\label{thm:transition-hierarchy}
\status{A}
For every fixed word $w$,
\begin{equation}\label{eq:transition-hierarchy}
 A_w(N)\ll_w\frac{N}{(\log N)^{\tau(w)}}.
\end{equation}
If $\tau(w)\ge2$, the sum of the reciprocals of its occurrence positions converges.
\end{theorem}

\begin{proof}
Each transition in $w$ forces, through \eqref{eq:xor}, a prime at a distinct fixed offset
from the starting point.  A fixed-tuple upper-bound sieve gives
\eqref{eq:transition-hierarchy}; a local congruence obstruction may make the actual count
eventually zero.  Partial summation proves reciprocal summability for
$\tau(w)\ge2$.
\end{proof}

\subsection{Prime-rooted walls and a reciprocal phase transition}

The usual integer-rooted local limit is dominated by long constant blocks.  Centering
instead at a uniformly selected prime exposes the transition itself.

\begin{theorem}[Prime-rooted domain-wall limit]\label{thm:prime-rooted-wall}
\status{A+E}
Fix $R\ge1$, choose $p_k\le X$ uniformly among the primes, and observe
\[
 (b_{p_k-R},\ldots,b_{p_k-1}\mid b_{p_k},\ldots,b_{p_k+R}).
\]
Its distribution converges, with cylinder error $O_R(1/\log X)$, to
\begin{equation}\label{eq:prime-rooted-wall}
 \frac12\delta_{0^\infty\mid1^\infty}
 +\frac12\delta_{1^\infty\mid0^\infty}.
\end{equation}
\end{theorem}

\begin{proof}
A centered block fails to be a single wall only if another prime lies within distance
$R$ of $p_k$.  Summing the fixed-shift prime-pair sieve bound over $1\le d\le R$
shows that there are $O_R(X/\log^2X)$ exceptional centers.  Division by $\pi(X)$
gives the stated error.  Odd- and even-indexed primes differ in number by at most one
and give the two opposite orientations.
\end{proof}

For a finite word $w$, let
\[
 H_w(X):=
 \sum_{\substack{n\le X\\b_nb_{n+1}\cdots b_{n+|w|-1}=w}}\frac1n.
\]

\begin{theorem}[Reciprocal three-phase law]\label{thm:reciprocal-three-phase}
\status{A+E}
There is a constant $\kappa_w$ in the one-transition case such that
\begin{equation}\label{eq:reciprocal-three-phase}
\begin{array}{ll}
 \tau(w)=0:&H_w(X)\asymp_w\log X,\\[1mm]
 \tau(w)=1:&H_w(X)=\dfrac12\log\log X+\kappa_w
                  +O_w(1/\log X),\\[2mm]
 \tau(w)\ge2:&H_w(\infty)-H_w(X)
                  =O_w((\log X)^{1-\tau(w)}).
\end{array}
\end{equation}
For $\tau(w)=1$ one has more precisely
\begin{equation}\label{eq:one-transition-count}
 A_w(X)=\frac12\pi(X)+O_w(X/\log^2X).
\end{equation}
\end{theorem}

\begin{proof}
If $\tau(w)=1$, its unique transition specifies one prime at a fixed offset and one
of the two prime-index parities.  These parities split the first $\pi(X)$ primes equally
up to one; an additional prime inside the fixed window accounts for
$O_w(X/\log^2X)$ exceptions.  This proves \eqref{eq:one-transition-count}, and
partial summation together with the Mertens theorem for reciprocal primes gives the
middle line of \eqref{eq:reciprocal-three-phase}.  For $\tau(w)\ge2$, apply partial
summation to \eqref{eq:transition-hierarchy}.  Finally, a constant word occurs on a
positive lower density of positions by \eqref{eq:parity-positive-density}, after
discarding the zero-density set lying within a fixed distance of a prime; the trivial
upper bound completes the first line.
\end{proof}

\subsection{Finite diffraction and its first high-pass scale}

Put
\[
 S_N(\theta):=\sum_{n=1}^N\sigma_ne^{in\theta},
 \qquad
 d\widehat\mu_N(\theta):=
 \frac{|S_N(\theta)|^2}{N}\frac{\dd\theta}{2\pi}.
\]
This is a probability measure on the circle.  Write
$d_{\rm ch}(e^{i\theta},1)=|e^{i\theta}-1|$.

\begin{proposition}[Exact prime count from finite diffraction]
\label{prop:diffraction-prime-count}
\status{E}
\begin{equation}\label{eq:diffraction-W2}
 W_{2,{\rm ch}}(\widehat\mu_N,\delta_0)^2
 =\frac{4\pi(N)+2}{N},
 \qquad
 \pi(N)=\frac N4W_{2,{\rm ch}}(\widehat\mu_N,\delta_0)^2-\frac12.
\end{equation}
Consequently the prime number theorem is equivalent to
\[
 \sqrt{\log N}\,W_{2,{\rm ch}}(\widehat\mu_N,\delta_0)\longrightarrow2.
\]
\end{proposition}

\begin{proof}
The only coupling with the point mass $\delta_0$ gives
\[
 W_{2,{\rm ch}}^2=2-2\Re\widehat{\widehat\mu_N}(1).
\]
Parseval and the adjacent sign change at every prime give
\[
 \widehat{\widehat\mu_N}(1)
 =\frac1N\sum_{n<N}\sigma_n\sigma_{n+1}
 =\frac{N-1-2\pi(N)}N,
\]
which is \eqref{eq:diffraction-W2}.
\end{proof}

\begin{theorem}[Renormalized high-pass diffraction]\label{thm:white-high-pass}
\status{A+E}
Define
\[
 d\widehat\nu_N(\theta)
 :=(\log N)\sin^2\!\frac\theta2\,d\widehat\mu_N(\theta).
\]
Then
\begin{equation}\label{eq:white-high-pass}
 \widehat\nu_N\Longrightarrow\frac{\dd\theta}{2\pi}.
\end{equation}
Thus the unfiltered diffraction converges to the Bragg mass $\delta_0$, whereas its
first nontrivial high-pass scale is Haar measure.
\end{theorem}

\begin{proof}
For fixed $h\in\mathbb Z$, put
\[
 r_N(h)=\frac1N\sum_{n=1}^{N-|h|}\sigma_n\sigma_{n+|h|}.
\]
By \cref{thm:fixed-lag},
$r_N(h)=1-2|h|/\log N+O_h(\log^{-2}N)$.  Hence
\[
 \widehat{\widehat\nu_N}(h)
 =\log N\left(\frac12r_N(h)-\frac14r_N(h-1)-\frac14r_N(h+1)\right).
\]
The discrete second difference of $|h|$ is $2$ at $h=0$ and zero elsewhere.
The displayed Fourier coefficient therefore tends to $1$ at $h=0$ and to $0$
otherwise.  Positivity and Fourier uniqueness on the circle prove
\eqref{eq:white-high-pass}.
\end{proof}

\begin{corollary}[Pure diffraction without Besicovitch phase coherence]
\label{cor:not-B2-AP}
\status{L+E}
With respect to the prefix F\o lner sequence $[1,N]$, the prime-parity sign comb has
pure-point diffraction $\delta_0$, but $(\sigma_n)$ is not Besicovitch
$B^2$-almost periodic.
\end{corollary}

\begin{proof}
The consistent-phase characterization of Besicovitch almost periodic diffraction
\cite{LenzSpindelerStrungaru} would force a $B^2$-almost periodic sequence with the
single diffraction atom of mass one to agree in $B^2$ with a constant phase.  This is
impossible here because \eqref{eq:parity-positive-density} gives positive lower density
to each of the values $+1$ and $-1$.
\end{proof}

\section{Correlation curvature and individual prime-gap thresholds}

\subsection{The discrete correlation Laplacian}

For a finitely supported sequence $u$ set $(\Delta u)_n=u_n-u_{n-1}$ and
$T_u(h)=\sum_nu_nu_{n+h}$.

\begin{proposition}[Correlation curvature]\label{prop:correlation-laplacian}
\status{E}
\begin{equation}\label{eq:correlation-laplacian}
 T_{\Delta u}(h)=2T_u(h)-T_u(h-1)-T_u(h+1).
\end{equation}
\end{proposition}

\begin{proof}
Expand $(1-S)$ in both factors and shift the summation indices.  For truncations of an
infinite sequence, the corresponding endpoint terms must be retained.
\end{proof}

Thus the negative discrete Laplacian of a dense cumulative correlation recovers the
correlation of its sparse signed difference.

\subsection{All-lag Green inversion and exact parity recovery}

For $X=p_{2m}$, put
\[
 C_X(h):=
 \sum_{\substack{i<j,\ p_j\le X\\p_j-p_i=h}}(-1)^{j-i},
 \qquad
 \mathcal T(x):=\sum_{n\le x}(-1)^{\pi(n)},
\]
and
\[
 D_X(h):=\sum_{n=1}^{X-1-h}b_nb_{n+h}\qquad(h\ge0),
\]
with $D_X(h)=0$ when the sum is empty.

\begin{theorem}[All-lag Green inversion]\label{thm:all-lag-green}
\status{E}
For every $h\ge1$,
\begin{equation}\label{eq:all-lag-curvature}
 C_X(h)=2D_X(h)-D_X(h-1)-D_X(h+1).
\end{equation}
Conversely, for every $h\ge0$,
\begin{equation}\label{eq:all-lag-green}
 D_X(h)=-\sum_{t>h}(t-h)C_X(t).
\end{equation}
In particular,
\begin{align}
 \#\{n\le X-1:\pi(n)\text{ is odd}\}
  &=-\sum_{h\ge1}hC_X(h),\label{eq:odd-count-green}\\
 \mathcal T(X-1)
  &=X-1+2\sum_{h\ge1}hC_X(h).\label{eq:T-green}
\end{align}
The correlation also obeys the exact sum rules
\begin{equation}\label{eq:CX-sum-rules}
 \sum_{h\ge1}C_X(h)=-m,
 \qquad
 \sum_{h\ge1}h^2C_X(h)
 =-\left(\sum_{h\ge1}hC_X(h)\right)^2.
\end{equation}
\end{theorem}

\begin{proof}
Extend $u_n=b_n\one_{1\le n\le X-1}$ by zero to $\mathbb Z$.  Since
$X=p_{2m}$, its difference is supported exactly on $p_1,\ldots,p_{2m}$ and
$(\Delta u)_{p_k}=(-1)^{k+1}$.  Hence $T_{\Delta u}(h)=C_X(h)$, so
\cref{prop:correlation-laplacian} gives \eqref{eq:all-lag-curvature}.  Twofold
summation, using the finite support of $D_X$, gives \eqref{eq:all-lag-green}; taking
$h=0$ yields \eqref{eq:odd-count-green} and \eqref{eq:T-green}.

Write $\epsilon_k=(-1)^{k+1}$.  Since $\sum_{k\le2m}\epsilon_k=0$,
\[
 \sum_hC_X(h)=\sum_{i<j}\epsilon_i\epsilon_j=-m,
 \qquad
 \sum_hhC_X(h)=\sum_k\epsilon_kp_k=-D_X(0).
\]
Expanding $(p_j-p_i)^2$ and again using $\sum\epsilon_k=0$ gives
\[
 \sum_hh^2C_X(h)
 =\sum_{i<j}\epsilon_i\epsilon_j(p_j-p_i)^2
 =-\left(\sum_k\epsilon_kp_k\right)^2,
\]
which completes \eqref{eq:CX-sum-rules}.
\end{proof}

Thus the open estimate $\mathcal T(x)=o(x)$ is, along even prime endpoints,
equivalent to the single signed-moment asymptotic
\begin{equation}\label{eq:T-moment-equivalence}
 \sum_{h\ge1}hC_{p_{2m}}(h)=-\frac{p_{2m}}2+o(p_{2m}).
\end{equation}
This identifies the required statistic but does not estimate it.

\subsection{Prime-rank phase tomography}

For $|z|=1$, set
\[
 u_z(n):=z^{\pi(n)},\qquad d_z(n):=u_z(n)-u_z(n-1).
\]
Then $d_z(p_k)=(z-1)z^{k-1}$ and $d_z(n)=0$ off the primes.  Let
\[
 G_{r,X}(h):=
 \#\{i:p_{i+r}-p_i=h,\ p_{i+r}\le X\}.
\]

\begin{proposition}[Exact phase decoder]\label{prop:phase-tomography}
\status{E}
For $z\ne1$,
\begin{equation}\label{eq:phase-polynomial}
 \mathcal F_{X,h}(z):=
 \frac1{|z-1|^2}\sum_{n\le X-h}\overline{d_z(n)}d_z(n+h)
 =\sum_{r=1}^{h}G_{r,X}(h)z^r.
\end{equation}
The value at $z=1$ is defined by polynomial continuation.  If $q>h$ and
$\omega=e^{2\pi i/q}$, then
\begin{equation}\label{eq:phase-DFT}
 G_{r,X}(h)=\frac1q\sum_{a=0}^{q-1}
 \omega^{-ar}\mathcal F_{X,h}(\omega^a).
\end{equation}
Thus all prime-index distances at a fixed additive distance are recovered exactly; in
particular, $G_{1,X}(h)$ is the number of consecutive prime gaps equal to $h$, with no
restriction $h<2H_*$.
\end{proposition}

\begin{proof}
Every nonzero summand on the left corresponds to a prime pair
$p_i,p_j=p_i+h$, and after division by $|z-1|^2$ its value is $z^{j-i}$.
Grouping by $r=j-i$ proves \eqref{eq:phase-polynomial}; finite Fourier inversion proves
\eqref{eq:phase-DFT}.
\end{proof}

The decoder reconstructs the counts $G_{r,X}(h)$, not the temporal order of the gaps.
Recovering ordered finite gap words requires $X$-increments or higher multivariate
correlations.

\subsection{An exact twin-prime curvature identity}

Put
\[
 R_N(h):=\sum_{n=1}^{N}\sigma_n\sigma_{n+h},
 \qquad
 \Pi_2(X):=\#\{p\le X-2:p,p+2\in\PP\}.
\]

\begin{theorem}[Twin-prime curvature]\label{thm:twin-curvature}
\status{E}
\begin{equation}\label{eq:twin-curvature}
 R_N(3)-2R_N(2)+R_N(1)
 =4\Pi_2(N+3)-2-2\chi_{N+3}.
\end{equation}
\end{theorem}

\begin{proof}
Use
$\sigma_n\sigma_{n+h}=\prod_{j=1}^{h}(1-2\chi_{n+j})$ for $h=1,2,3$.
In the displayed second difference, the interior linear terms telescope and the
nonadjacent quadratic term $4\chi_{n+1}\chi_{n+3}$ counts twin pairs.  The only adjacent
prime pair is $(2,3)$; collecting it and the two endpoints gives
\eqref{eq:twin-curvature}.
\end{proof}

Twin-prime infinitude is equivalent to unboundedness above of the left side.  The theorem
is an exact reformulation, not a new lower bound for prime pairs.

\subsection{Cumulative and individual signed correlations}

Recall the signed correlation $C_X(h)$ defined above, and put
\[
 L_{i,X}(H):=\#\{j>i:p_j\le X,\ p_j-p_i\le H\}.
\]

\begin{theorem}[Cumulative anti-correlation]\label{thm:cumulative-anti}
\status{E}
\begin{equation}\label{eq:cumulative-anti}
 \sum_{h=1}^{H}C_X(h)
 =-\#\{i:L_{i,X}(H)\text{ is odd}\}\le0.
\end{equation}
\end{theorem}

\begin{proof}
For fixed $i$, its contribution is
$\sum_{d=1}^{L_{i,X}(H)}(-1)^d$, equal to $-1$ for odd $L_{i,X}(H)$ and $0$ for even
$L_{i,X}(H)$.
\end{proof}

\begin{theorem}[Bounded-gap criterion]\label{thm:bounded-gap-criterion}
\status{E}
For fixed $H$, the following are equivalent:
\[
 g_i\le H\text{ for infinitely many }i,
 \qquad
 -\sum_{h\le H}C_X(h)\longrightarrow\infty.
\]
Consequently, in the extended integers,
\begin{equation}\label{eq:Hanti}
 \inf\left\{H:-\sum_{h\le H}C_X(h)\to\infty\right\}
 =\liminf_{n\to\infty}g_n.
\end{equation}
\end{theorem}

\begin{proof}
The divergence in \eqref{eq:cumulative-anti} immediately implies infinitely many
$g_i\le H$.  Conversely, if the right-hand side of \eqref{eq:cumulative-anti} did not
tend to infinity, then only finitely many of the limiting counts
$L_i(H):=\#\{j>i:p_j-p_i\le H\}$ could be odd: otherwise any prescribed finite set of
odd limiting counts would eventually appear among the $L_{i,X}(H)$.  Thus eventually
every $L_i(H)$ is even.  Whenever $g_i\le H$, we must then have $L_i(H)\ge2$, forcing
$g_{i+1}\le H$.  Induction would make every
sufficiently late gap at most $H$, implying $\pi(x)\gg x/H$, contrary to
$\pi(x)=o(x)$.  Taking the least integer threshold gives \eqref{eq:Hanti}.
\end{proof}

The cumulative identity can be sharpened at each individual lag.  Let
\[
 \Pi_h(X):=\#\{p:p+h\le X,\ p,p+h\in\PP\},
\]
and let $\Pi_{\{0,t,h\}}(X)$ count $p\le X-h$ for which
$p,p+t,p+h$ are all prime.

\begin{theorem}[Prime-pair main term]\label{thm:CX-pair-main}
\status{A+E}
For every fixed $h$,
\begin{equation}\label{eq:CX-exact-comparison}
 0\le C_X(h)+\Pi_h(X)
 \le2\sum_{1\le t<h}\Pi_{\{0,t,h\}}(X),
\end{equation}
and hence
\begin{equation}\label{eq:CX-asymptotic}
 C_X(h)=-\Pi_h(X)+O_h\!\left(\frac{X}{\log^3X}\right).
\end{equation}
In particular,
\begin{equation}\label{eq:CX2}
 C_X(2)=-\Pi_2(X)
\end{equation}
exactly.
\end{theorem}

\begin{proof}
A pair at distance $h$ contributes $-1$ when its prime-index difference is odd and $+1$
when that difference is even.  After adding the unweighted count $\Pi_h(X)$, the first
case contributes $0$ and the second $2$.  An even index difference is at least $2$, so
there is an intermediate prime $p+t$, giving \eqref{eq:CX-exact-comparison}.  A
fixed-triple Selberg upper-bound sieve gives
$\Pi_{\{0,t,h\}}(X)\ll_hX/\log^3X$.  For $h=2$, the only internal offset is $t=1$,
and no three integers $p,p+1,p+2$ are all prime; hence the correction term vanishes.
\end{proof}

\begin{theorem}[Single-lag recovery of the least recurrent gap]\label{thm:single-lag-gap}
\status{E}
Let
\[
 H_*:=\liminf_{n\to\infty}g_n.
\]
For every fixed $h<2H_*$,
\begin{equation}\label{eq:CX-consecutive-gap}
 C_X(h)=-\#\{n:p_{n+1}\le X,\ g_n=h\}+O_h(1).
\end{equation}
Therefore
\begin{equation}\label{eq:single-lag-threshold}
 H_*=\min\{h\ge1:C_X(h)\to-\infty\}.
\end{equation}
\end{theorem}

\begin{proof}
Because the gaps are integral, all sufficiently late gaps are at least $H_*$, while
$g_n=H_*$ occurs infinitely often.  A sufficiently late prime pair at distance
$h<2H_*$ cannot contain an intermediate prime, since then its distance would be the sum
of at least two gaps, each at least $H_*$.  This proves \eqref{eq:CX-consecutive-gap}.
For $h<H_*$ there are only finitely many such pairs, so $C_X(h)$ is eventually constant;
for $h=H_*$ the right side tends to $-\infty$.
\end{proof}

The published Polymath8b bound $H_*\le246$ \cite{Polymath8b} now gives the stronger
individual conclusion
\begin{equation}\label{eq:published-246}
 \text{for some even }h\le246,\qquad C_X(h)\longrightarrow-\infty.
\end{equation}
This is full divergence, rather than merely a divergent liminf.  It does not identify the
lag $h$ and does not imply the twin-prime conjecture $H_*=2$.

\section{Doubling orbits, Gauss orbits, and continuant ratios}

\subsection{Endpoint depth in the doubling orbit}

At the $k$-th prime time define
\begin{equation}\label{eq:xk-delta}
 x_k:=\left\{2^{p_k-1}(2C_{PP})\right\},\qquad
 d_k:=k\pmod2,\qquad
 \delta_k:=|x_k-d_k|.
\end{equation}
Thus $d_k\in\{0,1\}$ and $\delta_k$ is the distance from $x_k$ to the nearest endpoint
of the unit interval.

\begin{theorem}[Affine gap cocycle and binary decoder]\label{thm:binary-gap-decoder}
\status{E}
For every $k\ge1$,
\begin{align}
 x_k&=d_k(1-2^{-g_k})+2^{-g_k}x_{k+1},\label{eq:affine-gap-cocycle}\\
 2^{-(g_k+1)}&<\delta_k<2^{-g_k},\qquad
 g_k=\left\lfloor-\log_2\delta_k\right\rfloor.\label{eq:delta-decoder}
\end{align}
\end{theorem}

\begin{proof}
The binary expansion of $x_k$ begins with exactly $g_k$ copies of $d_k$, followed by
the opposite bit.  Summing this initial run gives \eqref{eq:affine-gap-cocycle}.  If
$d_k=0$, then $x_{k+1}\in(1/2,1)$; if $d_k=1$, then
$x_{k+1}\in(0,1/2)$.  The two strict bounds and the floor decoder follow.
\end{proof}

\subsection{The corrected Gauss-tail dictionary}

Let
\begin{equation}\label{eq:alphaP}
 \alpha_P=[0;2,g_1,g_2,\ldots]
\end{equation}
and write its convergent denominators as
\begin{equation}\label{eq:q-recurrence}
 q_{-1}=0,\quad q_0=1,\quad q_1=2,
 \qquad q_{n+1}=g_nq_n+q_{n-1}\quad(n\ge1).
\end{equation}
For the Gauss map $G(x)=\{1/x\}$, the correct indexing is
\begin{equation}\label{eq:zk-corrected}
 z_k:=G^k(\alpha_P)=[0;g_k,g_{k+1},\ldots]\qquad(k\ge1).
\end{equation}

\begin{theorem}[Three exact integer parts]\label{thm:three-decoders}
\status{E}
For every $k\ge1$,
\begin{equation}\label{eq:three-decoders}
 \left\lfloor\frac{q_{k+1}}{q_k}\right\rfloor
 =\left\lfloor\frac1{z_k}\right\rfloor
 =\left\lfloor-\log_2\delta_k\right\rfloor
 =g_k.
\end{equation}
In particular,
\begin{equation}\label{eq:cusp-distance-one}
 \left|\frac1{z_k}+\log_2\delta_k\right|<1.
\end{equation}
\end{theorem}

\begin{proof}
The recurrence gives
$q_{k+1}/q_k=g_k+q_{k-1}/q_k$ with the second term in $(0,1)$.  The continued-fraction
tail gives $1/z_k=g_k+z_{k+1}$ with $z_{k+1}\in(0,1)$.  The third equality is
\eqref{eq:delta-decoder}; the two real numbers in \eqref{eq:cusp-distance-one} lie in the
same unit interval.
\end{proof}

\begin{theorem}[Bounded-gap dictionary]\label{thm:bounded-gap-dictionary}
\status{E}
For a fixed integer $H\ge1$, the following are equivalent:
\begin{enumerate}[label=\textup{(\alph*)}]
 \item $-\sum_{h\le H}C_X(h)\to\infty$;
 \item $\delta_k>2^{-(H+1)}$ for infinitely many $k$;
 \item $z_k>1/(H+1)$ for infinitely many $k$;
 \item $\lfloor q_{k+1}/q_k\rfloor\le H$ for infinitely many $k$.
\end{enumerate}
\end{theorem}

\begin{proof}
By \cref{thm:bounded-gap-criterion}, (a) is equivalent to $g_k\le H$ infinitely often.
The equivalence with (b) is \eqref{eq:delta-decoder}; that with (d) is
\eqref{eq:three-decoders}.  Finally
$z_k=1/(g_k+z_{k+1})$, with $0<z_{k+1}<1$, makes $g_k\le H$ equivalent to
$z_k>1/(H+1)$.
\end{proof}

This theorem is a lossless coordinate translation of bounded gaps.  It does not improve
the numerical bounded-gap constant by itself.

\subsection{Density-one cusp concentration and recurrent excursions}

\begin{lemma}[Small consecutive gaps]\label{lem:small-gaps}
\status{A}
Uniformly for $2\le y\le\log X$,
\begin{equation}\label{eq:small-gaps}
 \#\{k:p_{k+1}\le X,\ g_k\le y\}
 \ll\frac{Xy}{\log^2X}.
\end{equation}
\end{lemma}

\begin{proof}
Enlarge the count of consecutive pairs to all prime pairs at even distances $h\le y$.
The Selberg pair sieve gives
\[
 \#\{p\le X:p,p+h\in\PP\}
 \ll \frac{X}{\log^2X}
 \prod_{\substack{\ell\mid h\\\ell>2}}\frac{\ell-1}{\ell-2}.
\]
Expanding the product as a squarefree divisor sum shows that its average over $h\le y$
is $O(1)$, yielding \eqref{eq:small-gaps}.
\end{proof}

For every fixed $H$,
\begin{equation}\label{eq:bounded-excursion-density}
 \#\{k\le K:g_k\le H\}=O_H(K/\log K).
\end{equation}
Nevertheless the set is infinite for $H=246$ by Polymath8b.  Thus
$g_k\to\infty$, $\delta_k\to0$, $z_k\to0$, and $q_{k+1}/q_k\to\infty$ each hold on
a set of indices of natural density one, but none holds pointwise.

Let
\[
 \nu_K:=\frac1K\sum_{k\le K}\delta_{\delta_k},
 \qquad
 \eta_K:=\frac1K\sum_{k\le K}\delta_{z_k}.
\]

\begin{theorem}[Quantitative cusp concentration]\label{thm:cusp-Wasserstein}
\status{A}
For the usual metric on $[0,1]$,
\begin{equation}\label{eq:cusp-Wasserstein}
 W_1(\nu_K,\delta_0)\ll\frac1{\log K},
 \qquad
 W_1(\eta_K,\delta_0)\ll\frac{\log\log K}{\log K}.
\end{equation}
\end{theorem}

\begin{proof}
Put $X=p_{K+1}\asymp K\log K$.  From \cref{lem:small-gaps} and the prime number
theorem,
\[
 \#\{k\le K:g_k\le y\}\ll K\frac{y}{\log K}
 \qquad(2\le y\le\log X).
\]
Summation over dyadic gap layers gives
\[
 \sum_{k\le K}2^{-g_k}\ll\frac K{\log K},
 \qquad
 \sum_{k\le K}\frac1{g_k}\ll
 \frac{K\log\log K}{\log K}.
\]
Indeed the first series has a summable weight on layers, while each of the
$O(\log\log K)$ dyadic layers up to $\log X$ contributes $O(K/\log K)$ to the second;
the tail $g_k>\log X$ contributes $O(K/\log X)$.  Now
$\delta_k<2^{-g_k}$ and $z_k<1/g_k$.  The $W_1$ distance from an empirical measure on
$[0,1]$ to $\delta_0$ is the mean of its coordinates, proving the result.
\end{proof}

\subsection{A dimension bound for the Gauss-tail closure}

Weak concentration at the cusp does not by itself control the closure of the exceptional
tails.  The residue graph modulo $6$ supplies such a geometric bound.

\begin{theorem}[Thin Gauss-tail closure]\label{thm:gauss-closure-dimension}
\status{A+E}
For
\[
 Z:=\overline{\{z_k:k\ge2\}}\subset[0,1],
\]
one has
\begin{equation}\label{eq:gauss-closure-dimension}
 \dim_H Z\le\frac23.
\end{equation}
\end{theorem}

\begin{proof}
For primes larger than $3$, the residue of $p_k$ modulo $6$ is $1$ or $5$.
The inverse branch associated with a gap $g$ is
$\phi_g(x)=1/(g+x)$, with $|\phi_g'(x)|\le g^{-2}$.  At Hausdorff exponent
$s=2/3$, the two-state content matrix is dominated by
\[
 M=\begin{pmatrix}A&B\\ C&A\end{pmatrix},
\]
where
\[
 A=\sum_{m\ge1}(6m)^{-4/3},\quad
 B=\sum_{m\ge0}(6m+4)^{-4/3},\quad
 C=\sum_{m\ge0}(6m+2)^{-4/3}.
\]
The integral test gives
\[
 A\le\frac4{6^{4/3}},\qquad
 B\le4^{-4/3}+\frac12\,4^{-1/3},\qquad
 C\le2^{-4/3}+\frac12\,2^{-1/3}.
\]
Consequently
\[
 \rho(M)=A+\sqrt{BC}<0.980<1.
\]
Iterating the graph-directed cylinder cover therefore makes its total $2/3$-content
tend to zero.  The finitely many initial tails do not affect Hausdorff dimension, proving
\eqref{eq:gauss-closure-dimension}.
\end{proof}

\section{Prime-gap continuants and lacunary support series}

\subsection{OEIS continuants and logarithmic growth}

Let
\[
 \beta:=[1;g_1,g_2,\ldots]
\]
and let $A_n/B_n=[1;g_1,\ldots,g_n]$ with offset $n=0$.  Then
\begin{equation}\label{eq:OEIS-CF-dictionary}
 \alpha_P=\frac1{1+\beta},\qquad q_{n+1}=A_n+B_n.
\end{equation}
Thus for $Q_m:=q_{2m}$,
\begin{equation}\label{eq:Q-OEIS}
 Q_m=A_{2m-1}+B_{2m-1};
\end{equation}
$Q_m$ should not be identified with the denominator sequence $B_n$ itself.  The
sequences $A_n,B_n$ are OEIS A109139 and A109140 with their stated offset.

For positive partial quotients, the continuant recurrence gives the elementary sandwich
\[
 \prod_{j\le n}a_j\le K_n(a_1,\ldots,a_n)
 \le\prod_{j\le n}(a_j+1).
\]
Here the exact recurrence gives a slightly sharper form.

\begin{theorem}[Full-gap denominator growth]\label{thm:q-growth}
\status{A}
As $K\to\infty$,
\begin{equation}\label{eq:q-growth-precise}
 0\le
 \log q_{K+1}-\log2-\sum_{k\le K}\log g_k
 \le\sum_{k\le K}\log\left(1+\frac1{g_k}\right)
 \ll\frac{K\log\log K}{\log K}.
\end{equation}
Consequently
\begin{equation}\label{eq:q-growth}
 \log q_n=n\log\log n+O(n),
 \qquad
 \log Q_m=2m\log\log m+O(m).
\end{equation}
\end{theorem}

\begin{proof}
The continuant inequalities give the first two bounds in
\eqref{eq:q-growth-precise}, and the layer estimate in the proof of
\cref{thm:cusp-Wasserstein} bounds the last sum.

Let $X=p_{K+1}\asymp K\log K$.  Integrating \cref{lem:small-gaps} over the levels
$g_k<\log X$ yields
\[
 \sum_{k\le K}(\log\log X-\log g_k)_+=O(K).
\]
Hence $\sum_{k\le K}\log g_k\ge K\log\log X-O(K)$.  Jensen's inequality and
$\sum_{k\le K}g_k=p_{K+1}-2$ give the reverse bound
\[
 \sum_{k\le K}\log g_k
 \le K\log\frac{p_{K+1}-2}{K}
 =K\log\log K+O(K).
\]
Combine these estimates with \eqref{eq:q-growth-precise}.
\end{proof}

\subsection{Half gaps and irrationality exponents}

Write the odd primes as $3=r_1<r_2<\cdots$ and set
\[
 h_n:=\frac{r_{n+1}-r_n}{2},\qquad
 H_{-1}=0,\quad H_0=1,\quad
 H_n=h_nH_{n-1}+H_{n-2}.
\]

\begin{theorem}[Half-gap growth and three irrationality exponents]
\label{thm:half-gap-growth}
\status{A}
\begin{equation}\label{eq:H-growth}
 \log H_n=n\log\log n+O(n),
\end{equation}
and
\begin{equation}\label{eq:mu-two}
 \mu(\alpha_P)=\mu(\beta)
 =\mu([0;h_1,h_2,\ldots])=2.
\end{equation}
\end{theorem}

\begin{proof}
Replacing $g_n$ by $h_n=g_{n+1}/2$ changes the relevant logarithmic sum by $O(n)$, so
the proof of \cref{thm:q-growth} gives \eqref{eq:H-growth}.  For an infinite simple
continued fraction,
\[
 \mu([a_0;a_1,a_2,\ldots])
 =2+\limsup_n\frac{\log a_{n+1}}{\log q_n}.
\]
Now $\log g_n,\log h_n=O(\log n)$, whereas the denominator logarithms are
$n\log\log n+O(n)$.  The extra term is zero.  The fractional-linear relation
$\alpha_P=1/(1+\beta)$ preserves irrationality exponent.
\end{proof}

The exponent $2$ does not decide algebraicity: algebraic irrationals and many
transcendental numbers have the same irrationality exponent.

\subsection{Congruence rigidity}

\begin{theorem}[Full-gap congruence and lacunarity]\label{thm:Q-congruence}
\status{E}
For every $m\ge1$,
\begin{equation}\label{eq:Q-mod4-gap}
 Q_m\equiv3\pmod4,\qquad Q_{m+1}>5Q_m.
\end{equation}
\end{theorem}

\begin{proof}
All $g_k$ with $k\ge2$ are even.  Reducing \eqref{eq:q-recurrence} modulo $2$ and
$4$ gives $q_{2m-1}$ even and
$q_{2m}\equiv q_{2m-2}\equiv3\pmod4$.  Also
$q_{2m+1}>2q_{2m}$ and
$q_{2m+2}>2q_{2m+1}+q_{2m}>5q_{2m}$.
\end{proof}

\begin{theorem}[Exclusion of the factor $3$]\label{thm:H-mod3}
\status{E}
For every $n\ge1$,
\begin{equation}\label{eq:H-mod3}
 H_n\equiv-r_{n+1}H_{n-1}\pmod3,\qquad 3\nmid H_n.
\end{equation}
\end{theorem}

\begin{proof}
The case $n=1$ is direct.  Since $2^{-1}\equiv-1\pmod3$,
$h_n\equiv r_n-r_{n+1}\pmod3$.  For $n\ge2$, substitute the induction hypothesis
$H_{n-1}\equiv-r_nH_{n-2}$ and use $r_n^2\equiv1\pmod3$.  Iteration gives
\[
 H_n\equiv(-1)^n\prod_{j=2}^{n+1}r_j\pmod3,
\]
whose factors are all units modulo $3$.
\end{proof}

\subsection{Natural boundaries and arithmetic values}

Define
\begin{equation}\label{eq:FQ-FH}
 F_Q(z):=\sum_{m\ge1}z^{Q_m},
 \qquad
 F_H(z):=\sum_{n\ge0}z^{H_n}
          =2z+\sum_{n\ge2}z^{H_n}.
\end{equation}

\begin{theorem}[Support-series natural boundaries]\label{thm:FQ-FH-natural}
\status{A}
Both $F_Q$ and $F_H$ have the unit circle as a natural boundary.  They are
transcendental over $\mathbb C(z)$ and are not D-finite.
\end{theorem}

\begin{proof}
The inequality $Q_{m+1}/Q_m>5$ permits the Hadamard gap theorem for $F_Q$.  For the
half-gap sequence,
\[
 H_{n+1}-H_n=(h_{n+1}-1)H_n+H_{n-1}\ge H_{n-1}\to\infty.
\]
Apart from the initial coefficient $[z]F_H(z)=2$, its coefficients are $0$ or $1$.
They cannot be eventually periodic because the support gaps tend to infinity while the
support is infinite.  The radius is $1$, so P\'olya--Carlson gives the natural boundary.
The transcendence and non-D-finiteness consequences are as in
\cref{thm:AB-natural-boundary}.
\end{proof}

The congruence $Q_m\equiv3\pmod4$ also gives the exact rotation law
\begin{equation}\label{eq:FQ-rotation}
 F_Q(iz)=-iF_Q(z),\qquad F_Q(z)=z^3G(z^4)
\end{equation}
for a suitable power series $G$.

\begin{theorem}[Transcendental algebraic specializations]\label{thm:FQ-FH-algebraic}
\status{L}
For every algebraic number $\alpha$ with $0<|\alpha|<1$, both
$F_Q(\alpha)$ and $F_H(\alpha)$ are transcendental.
\end{theorem}

\begin{proof}
For $F_Q$, apply Corollary~5 of Corvaja--Zannier \cite{CorvajaZannier} to the Hadamard
sequence $Q_m$.  For $F_H$, use their Corollary~4: the exponent differences tend to
infinity, while
\[
 \limsup_n\frac{H_{n+1}}{H_n}
 \ge\limsup_n h_{n+1}=\infty
\]
because prime gaps are unbounded.  The finite initial coefficient $2z$ is harmless.
\end{proof}

\begin{lemma}[Lacunary rational specializations]\label{lem:Liouville-lacunary}
\status{E}
Let $e_1<e_2<\cdots$ be integers such that
\[
 e_{j+1}-e_j\to\infty,\qquad
 \limsup_j\frac{e_{j+1}}{e_j}=\infty.
\]
For every rational $r$ with $0<|r|<1$, the number $\sum_jr^{e_j}$ is Liouville.
\end{lemma}

\begin{proof}
Write $r=a/b$ in lowest terms.  The $J$-th partial sum has denominator dividing
$b^{e_J}$, while
\[
 \left|\sum_{j>J}r^{e_j}\right|
 \le\frac{|r|^{e_{J+1}}}{1-|r|}
\]
for large $J$; when $r<0$, the first omitted term dominates the remaining tail once the
support gaps are large.  Along a subsequence on which $e_{J+1}/e_J$ is arbitrarily large,
the error is smaller than every prescribed power of $b^{-e_J}$.  The tail is nonzero, and
the standard rational lower bound also rules out a rational limit.  Thus there are
infinitely many rational approximants of every exponent.
\end{proof}

\begin{corollary}[Liouville values of the two continuant series]
\label{cor:FQ-FH-Liouville}
\status{E}
For every rational $r$ with $0<|r|<1$, both $F_Q(r)$ and $F_H(r)$ are Liouville
numbers.
\end{corollary}

\begin{proof}
For $Q_m$, the ratios exceed $5$ and become arbitrarily large because
$Q_{m+1}/Q_m>g_{2m}g_{2m+1}$.  For $H_n$,
$H_{n+1}/H_n>h_{n+1}$ has infinite limsup and the exponent differences tend to
infinity.  Apply \cref{lem:Liouville-lacunary}.
\end{proof}

\section{Gilbreath obstruction languages at fixed depth}

This section concerns exact finite-state descriptions of obstruction languages, not a
proof of Gilbreath's conjecture.  Normalize the gaps by
\begin{equation}\label{eq:Gilbreath-c}
 c_j:=\frac{p_{j+2}-p_{j+1}}2-1,
 \qquad (\Dabs v)_j:=|v_{j+1}-v_j|,
 \qquad C_{r,j}:=(\Dabs^rc)_j.
\end{equation}
The equivalence with the classical Gilbreath formulation and the relevance of long zero
blocks and shallow $\{0,d\}$ blocks are due to Chase, Hunter and Tao
\cite{CHT}.  We develop the inverse language and its exact finite-state compression.

\subsection{Shallow arithmetic barriers}

\begin{proposition}[Depth-zero barriers]\label{prop:depth-zero}
\status{E}
The only two consecutive zero entries of $c$ are $c_1=c_2=0$.  More generally, if an
odd prime $\ell$ divides $d$ and $p_{j+1}>\ell$, then a $\{0,d\}$ block in $c$ has at
most $\ell-2$ entries.
\end{proposition}

\begin{proof}
Two consecutive zeros correspond to $p,p+2,p+4$, so reduction modulo $3$ leaves only
$3,5,7$.  In a $\{0,d\}$ block, every actual gap is
$2(c_t+1)\equiv2\pmod\ell$.  If there were $\ell-1$ gaps, their $\ell$ endpoints would
visit every residue modulo $\ell$, forcing one endpoint to be a multiple of $\ell$ larger
than $\ell$.
\end{proof}

\begin{proposition}[A depth-one primorial barrier]\label{prop:depth-one}
\status{A}
Suppose $L$ consecutive entries of $\Dabs c$ vanish, so that $L+1$ consecutive entries
of $c$ have a common value $u\le H$.  If the first prime endpoint exceeds $L+2$, then
\begin{equation}\label{eq:primorial-barrier}
 \prod_{\substack{3\le\ell\le L+2\\\ell\ \mathrm{prime}}}\ell\mid u+1.
\end{equation}
Consequently $L=O(\log(H+2))$.
\end{proposition}

\begin{proof}
The $L+2$ prime endpoints form an arithmetic progression of common difference
$2(u+1)$.  For each odd prime $\ell\le L+2$, a nonzero common difference modulo
$\ell$ would make the endpoints hit residue $0$.  Hence every such $\ell$ divides
$u+1$.  The prime number theorem for the Chebyshev function gives the logarithmic bound.
\end{proof}

\subsection{Affine inverse language and parity recurrence}

Fix a bottom word $V^{(r)}$.  A one-row inverse is specified by a starting value
$C_s=V^{(s)}_0$ and signs $\eps_t^{(s)}\in\{-1,1\}$ through
\begin{equation}\label{eq:inverse-row}
 V^{(s)}_{t+1}=V^{(s)}_t+\eps_t^{(s)}V^{(s+1)}_t,
 \qquad V^{(s)}_t\ge0.
\end{equation}

\begin{theorem}[Affine polyhedral inverse decomposition]
\label{thm:affine-inverse}
\status{E}
Fix $r,K,d$ and a bottom word $V^{(r)}=d\eta$ with
$\eta\in\{0,1\}^K$.  For each full sign pattern $\eps$ there are integer matrices
$A_\eps,B_\eps$ such that
\begin{equation}\label{eq:affine-inverse}
 V^{(0)}=A_\eps C+dB_\eps\eta.
\end{equation}
The nonnegativity constraints are rational linear inequalities in $(C,\eta)$.
Accordingly the integral top-row preimages form a finite union, over sign patterns, of
affine images of lattice points in rational polyhedra.
\end{theorem}

\begin{proof}
Iterate \eqref{eq:inverse-row} upward.  For a fixed sign pattern every entry is affine in
the row-start vector $C$ and in $d\eta$; all remaining restrictions are precisely the
nonnegativity inequalities.
\end{proof}

Modulo $2$ the signs disappear:
\begin{equation}\label{eq:parity-D-r}
 (\Dabs^rc)_j\equiv\sum_{t=0}^{r}\binom rt c_{j+t}\pmod2.
\end{equation}
If $d$ is even and $\Dabs^rc\in\{0,d\}$, then $(1+E)^rc=0$ over $\Ftwo$.  With
$P=2^{\lceil\log_2r\rceil}$, the Frobenius identity gives
$c_{j+P}\equiv c_j\pmod2$ at all interior indices where the recurrence is defined.
The order-$r$ recurrence is determined by its first $r$ bits, so there are at most
$2^r$ parity seeds, independently of the block length.

\subsection{Wheel coordinates and syndetic cusp escape}

For a height bound $H$, put $B=H+1$ and
\[
 u_0=0,\qquad u_m=\sum_{t<m}(c_t+1),\qquad
 1\le u_{m+1}-u_m\le B.
\]
Let $S$ be a finite set of odd primes, $Q=\prod_{\ell\in S}\ell$, and let
$\Delta(Q)$ be the largest cyclic gap between reduced residues modulo $Q$.

\begin{theorem}[Exact wheel criterion]\label{thm:wheel-criterion}
\status{E}
An infinite increasing wheel-survivor path with increments in $\{1,\ldots,B\}$ exists
if and only if
\begin{equation}\label{eq:wheel-criterion}
 \Delta(Q)\le B.
\end{equation}
\end{theorem}

\begin{proof}
For each $\ell\in S$ the primality constraint forbids one residue.  The Chinese
remainder theorem aligns these residues to a single translate
$r+(\mathbb Z/Q\mathbb Z)^\times$.  If every cyclic gap is at most $B$, traverse
successive allowed residues periodically.  Conversely an infinite path must repeatedly
cross every cyclic gap, so a gap larger than $B$ is impossible.
\end{proof}

Known lower bounds for long sifted intervals, in particular those of
Ford--Green--Konyagin--Maynard--Tao \cite{FGKMT}, imply that for every sufficiently large
$B$ one may choose a finite wheel for which no infinite height-$H$ survivor exists and
whose finite-state path length is bounded by
\begin{equation}\label{eq:wheel-length}
 L(B)\le
 \exp\!\left(
 O\!\left(\frac{B\log_2B}{\log B\,\log_3B}\right)
 \right)=\exp(o(B)).
\end{equation}
Here $\log_j$ denotes the $j$-fold iterated logarithm.
The harmless factor $2$ between integer gaps and the normalized half-gap coordinate has
been absorbed into $B$.  We do not use the stronger, previously uncited claim for a
particular primorial wheel.

\begin{corollary}[Syndetic cusp escape]\label{cor:syndetic-cusp}
\status{L+E}
For every fixed $B$ there is $L(B)<\infty$ such that, sufficiently far out, every block
of $L(B)$ consecutive prime gaps contains some $g_k>2B$.  At the same index,
\begin{equation}\label{eq:syndetic-cusp}
 \delta_k<2^{-2B},\qquad z_k<\frac1{2B}.
\end{equation}
For large $B$, $L(B)$ may be chosen as in \eqref{eq:wheel-length}.
\end{corollary}

\begin{proof}
A block with all $g_k\le2B$ gives a normalized survivor with increments at most $B$.
After the finitely many wheel primes, actual prime endpoints avoid all forbidden residues.
The acyclic finite wheel therefore bounds the path length.  Apply
\eqref{eq:delta-decoder} and $z_k<1/g_k$.
\end{proof}

\subsection{Anchored congruence states and transfer matrices}

Fix $r,H$ and a finite wheel $S$.  Let the current right difference boundary be
$E=(e_0,\ldots,e_{r-1})$.  Appending $a\in\{0,\ldots,H\}$ gives
\begin{equation}\label{eq:boundary-update}
 f_0=a,\qquad f_k=|f_{k-1}-e_{k-1}|\quad(1\le k\le r).
\end{equation}
The new boundary is $(f_0,\ldots,f_{r-1})$ and the new bottom label is $f_r$.
Fix now an actual starting prime $P$.  For every $\ell\in S$, let
\begin{equation}\label{eq:anchored-residue-state}
 x_\ell:=P+\text{current actual offset}\pmod\ell.
\end{equation}
Appending $a$ gives
\begin{equation}\label{eq:anchored-residue-update}
 x'_\ell=x_\ell+2(a+1)\pmod\ell,
 \qquad x'_\ell\ne0.
\end{equation}

\begin{theorem}[Anchored label-preserving right congruence]
\label{thm:bisimulation}
\status{E}
For fixed $(r,H,S,P)$, the history map
\[
 \text{history}\longmapsto\bigl(E,(x_\ell)_{\ell\in S}\bigr)
\]
has a label-preserving right-congruence kernel on the set of already wheel-admissible
histories.  Equal states have, for every append value, the same bottom label, the same
wheel admissibility, the same next state, and hence identical admissible labelled futures.
\end{theorem}

\begin{proof}
The bottom label and the next difference boundary are determined by
\eqref{eq:boundary-update}; the anchored congruence decision and next residues are
determined by \eqref{eq:anchored-residue-update}.  Induction on continuation length
proves the assertion.
\end{proof}

\begin{remark}[Why the endpoint residue is necessary]
\label{rem:bisimulation-counterexample}
The relative-mask state used in Version~2.0 is not a right congruence for a fixed actual
start.  Take $r=1$, $H=2$, $S=\{5\}$, and $P\equiv1\pmod5$.  The histories
\[
 h_1=(0,2,1,2),\qquad h_2=(2,0,0,2)
\]
have actual-offset trajectories modulo $5$
\[
 (0,2,3,2,3),\qquad(0,1,3,0,1),
\]
respectively.  Both have $E=(2)$ and the same translated visited mask
$C_5=\{0,2,4\}$.  Appending $a=2$ gives the same bottom label $0$, but the first
endpoint is $1+4\equiv0\pmod5$ and is forbidden, whereas the second is
$1+2\equiv3\pmod5$ and is allowed.  Relative masks remain appropriate for the
different, existential question whether some starting translate avoids all visited
residues; they cannot replace \eqref{eq:anchored-residue-state} in the anchored
language.
\end{remark}

The smallest wheel already gives a uniform entropy saving.  Let
$n_j=\#\{0\le a\le H:a\equiv j\pmod3\}$.  On the two nonzero endpoint residues
modulo $3$, the anchored transition matrix is
\begin{equation}\label{eq:mod3-wheel-matrix}
 A_{3,H}=\begin{pmatrix}n_2&n_1\\n_0&n_2\end{pmatrix},
 \qquad
 \rho(A_{3,H})=n_2+\sqrt{n_0n_1}\le\frac{2(H+1)}3.
\end{equation}
Thus the number of length-$K$ height words surviving even this single congruence is at
most $\exp[-(\log(3/2)+o(1))K]$ times the full $(H+1)^K$ word count.  This is a
uniform congruence saving, not a proof that every surviving word is realized by primes.

For $d\ge1$ define
\begin{equation}\label{eq:N-rd}
 N_{r,d}(K,H):=
 \#\{c\in\{0,\ldots,H\}^{K+r}:\Dabs^rc\in\{0,d\}^{K}\}.
\end{equation}

\begin{theorem}[Finite transfer representation]\label{thm:transfer}
\status{E}
There exist a finite nonnegative matrix $Q_{r,H,d}$ and an initial multiplicity vector
$\mathbf m$ such that
\begin{equation}\label{eq:transfer}
 N_{r,d}(K,H)=\mathbf m^{\!T}Q_{r,H,d}^{K}\mathbf1.
\end{equation}
Moreover
\begin{equation}\label{eq:transfer-branch}
 N_{r,d}(K,H)\le(H+1)^r(3\cdot2^{r-1})^K.
\end{equation}
\end{theorem}

\begin{proof}
The right boundary in \eqref{eq:boundary-update} has finitely many possibilities and
append is a Markov transition.  In reverse, a prescribed bottom label $0$ has at most
$2^{r-1}$ sign choices and a label $d$ at most $2^r$, giving the per-step branching bound
$3\cdot2^{r-1}$.
\end{proof}

Define the fixed-height entropy by
\[
 h_{r,d}(H):=\limsup_{K\to\infty}\frac1K\log N_{r,d}(K,H).
\]
For $r=1$, direct path-graph diagonalization gives
\[
 \rho(Q_{1,H,d})=1+2\cos\frac{\pi}{\lfloor H/d\rfloor+2}<3.
\]
For the next depth one obtains the explicit domination
\begin{equation}\label{eq:r2-entropy}
 \rho(Q_{2,H,d})
 \le2+4\cos\frac{\pi}{\lfloor H/d\rfloor+2}<6,
\end{equation}
and hence
\[
 h_{2,d}(H)\le
 \log\!\left(2+4\cos\frac{\pi}{\lfloor H/d\rfloor+2}\right).
\]
Indeed the dominating gap-state matrix has entries
\[
 W_{y,t}=m(t)\one_{|t-y|\in\{0,d\}},
 \qquad m(0)=1,\quad m(t)=2\ (t>0),
\]
and each residue component is a finite path with a self-loop.  Its adjacency eigenvalue is
\[
 2\cos\!\left(\frac{\pi}{\lfloor H/d\rfloor+2}\right),
\]
and the multiplicity weight produces \eqref{eq:r2-entropy}.

These finite-state results count and compress potential obstruction words.  They do not
prove that every actual prime input avoids them, and therefore do not settle Gilbreath's
conjecture.

\section{The dyadic numerator cocycle}

\subsection{Exact recurrence}

Let
\begin{equation}\label{eq:Cn-Nn}
 C_n:=\sum_{k=1}^{n}(-1)^{k+1}2^{-p_k}
 =\frac{N_n}{2^{p_n}}.
\end{equation}
The numerator sequence $(N_n)$ is recorded as OEIS A122150.

\begin{theorem}[Dyadic numerator recurrence]\label{thm:N-recurrence}
\status{E}
\begin{equation}\label{eq:N-recurrence}
 N_1=1,\qquad
 N_n=2^{g_{n-1}}N_{n-1}+(-1)^{n+1}\quad(n\ge2).
\end{equation}
Every $N_n$ is positive and odd, so $2^{p_n}$ is the exact reduced denominator in
\eqref{eq:Cn-Nn}.
\end{theorem}

\begin{proof}
Raise $C_{n-1}$ to denominator $2^{p_n}$ and add the final term.  The first summand in
\eqref{eq:N-recurrence} is even and the second odd.  Positivity also follows from the
alternating series with strictly decreasing term magnitudes.
\end{proof}

\subsection{Autonomous and indexed decoding; non-holonomicity}

For odd $x>1$, put
\[
 \chi_4(x):=(-1)^{(x-1)/2},
 \qquad
 \mathscr D_N(x):=\oddpart\!\bigl(x-\chi_4(x)\bigr),
\]
where $\oddpart(m)=m/2^{v_2(m)}$.

\begin{theorem}[Autonomous numerator decoder]
\label{thm:autonomous-N-decoder}
\status{E}
For every $n\ge3$,
\begin{equation}\label{eq:autonomous-N-decoder}
 \mathscr D_N(N_n)=N_{n-1},
 \qquad
 g_{n-1}=v_2(N_n^2-1)-1.
\end{equation}
Moreover,
\begin{equation}\label{eq:arch-2adic-resonance}
 g_{n-1}=\lfloor\log_2N_n\rfloor-\lfloor\log_2N_{n-1}\rfloor.
\end{equation}
Thus a single integer $N_n$, together with the fixed seed
$(p_1,p_2)=(2,3)$, determines $n$, all preceding gaps, and all primes
$p_1,\ldots,p_n$.
\end{theorem}

\begin{proof}
Write $g=g_{n-1}$, $M=N_{n-1}$, and $\varepsilon=(-1)^{n+1}$.  Since
$g\ge2$ for $n\ge3$, the recurrence $N_n=2^gM+\varepsilon$ gives
$N_n\equiv\varepsilon\pmod4$, hence $\chi_4(N_n)=\varepsilon$.  This proves the
first identity.  Also
\[
 N_n^2-1=2^{g+1}M\bigl(2^{g-1}M+\varepsilon\bigr),
\]
whose final factor is odd, proving the valuation formula.

If $2^d\le M<2^{d+1}$, then
$2^{d+g}\le2^gM+\varepsilon<2^{d+g+1}$.  In the only marginal case
$\varepsilon=-1$, the odd number $M>1$ is strictly larger than $2^d$.
This proves \eqref{eq:arch-2adic-resonance}.  Iterating $\mathscr D_N$ terminates at
$N_2=1$; the number of iterations determines $n$, while the displayed valuation
recovers the gaps and $g_1=1$ is fixed.
\end{proof}

Let $\epsilon_n=n\bmod2$, and let $a^m$ denote a run of $m$ copies of the binary
digit $a$.  The stronger finite-word form is
\begin{equation}\label{eq:N-binary-hologram}
 (N_n-\epsilon_n)_2
 =1^{g_1}0^{g_2}1^{g_3}0^{g_4}\cdots
 \bigl((n-1)\bmod2\bigr)^{g_{n-1}}\qquad(n\ge2).
\end{equation}
Indeed, an odd $n$ appends $g_{n-1}$ zeroes in the recurrence, while an even $n$
uses
$N_n=2^{g_{n-1}}(N_{n-1}-1)+(2^{g_{n-1}}-1)$ and appends that many ones.

\begin{theorem}[Non-holonomicity of OEIS A122150]
\label{thm:N-not-P-recursive}
\status{A+E}
The sequence $(N_n)$ is not $P$-recursive.  Equivalently, its ordinary generating
series
\[
 \mathcal N(t):=\sum_{n\ge1}N_nt^n
\]
is not $D$-finite as a formal power series.
\end{theorem}

\begin{proof}
The recurrence gives
\[
 \frac{N_{n+1}}{N_n}\ge2^{g_n}-1.
\]
The large-gap theorem of Ford--Green--Konyagin--Maynard--Tao \cite{FGKMT} implies
$\limsup g_n/\log n=\infty$.  Hence this ratio exceeds every fixed power of $n$
along a subsequence.

If the positive increasing sequence $(N_n)$ satisfied
$\sum_{j=0}^{d}P_j(n)N_{n+j}=0$ with $P_d\not\equiv0$, then for all large $n$
\[
 \frac{N_{n+d}}{N_{n+d-1}}
 \le\sum_{j<d}\left|\frac{P_j(n)}{P_d(n)}\right|
       \frac{N_{n+j}}{N_{n+d-1}}
 \le n^{O(1)},
\]
because the sequence is increasing and rational functions have polynomial growth.
This contradicts the subsequence above.  The equivalence with formal $D$-finiteness is
standard \cite{Stanley}.
\end{proof}

\begin{theorem}[$2$-adic gap decoder]\label{thm:2adic-decoder}
\status{E}
For every $n\ge2$,
\begin{equation}\label{eq:2adic-decoder}
 g_{n-1}=v_2\!\left(N_n-(-1)^{n+1}\right).
\end{equation}
Thus the numerator sequence alone recovers every prime:
\begin{equation}\label{eq:prime-from-N}
 p_n=2+\sum_{j=2}^{n}
 v_2\!\left(N_j-(-1)^{j+1}\right).
\end{equation}
\end{theorem}

\begin{proof}
Subtract the final sign in \eqref{eq:N-recurrence}.  Since $N_{n-1}$ is odd, the exact
$2$-adic valuation is $g_{n-1}$.  Sum the recovered gaps.
\end{proof}

\begin{corollary}[A modulo-$24$ twin-gap sensor]\label{cor:mod24-sensor}
\status{E}
For $n\ge3$,
\begin{equation}\label{eq:mod24-sensor}
 g_{n-1}=2\quad\Longleftrightarrow\quad
 \begin{cases}
  N_n\equiv5\pmod{24},&n\text{ odd},\\
  N_n\equiv19\pmod{24},&n\text{ even}.
 \end{cases}
\end{equation}
\end{corollary}

\begin{proof}
For $n\ge3$, $g_{n-1}$ is even.  Modulo $3$, $2^{g_{n-1}}\equiv1$, and
\eqref{eq:N-recurrence} gives $N_n\equiv2$ for odd $n$ and $N_n\equiv1$ for even
$n$.  Modulo $8$, $2^g\equiv4$ when $g=2$ and $0$ when $g\ge4$.  Combine the two
conditions by the Chinese remainder theorem.
\end{proof}

The sensor is exact but, like \cref{thm:twin-curvature}, does not provide a lower bound
for the number of twin gaps.

\subsection{Gap-polynomial resonance and fixed-lag separation}

For $1\le m<n$ put
\begin{equation}\label{eq:Rmn}
 R_{m,n}:=\sum_{k=m+1}^{n}(-1)^{k+1}2^{p_n-p_k}.
\end{equation}
Its final summand is odd and all preceding summands are even, so $R_{m,n}$ is odd and
nonzero.

\begin{theorem}[Gap-polynomial resonance]\label{thm:gcd-resonance}
\status{E}
\begin{equation}\label{eq:gcd-resonance}
 N_n=2^{p_n-p_m}N_m+R_{m,n},\qquad
 \gcd(N_n,N_m)=\gcd(N_m,R_{m,n}).
\end{equation}
If $n=m+2$ and an odd prime $\ell$ divides $N_m$, then
\begin{equation}\label{eq:order-resonance}
 \ell\mid N_{m+2}\quad\Longleftrightarrow\quad
 \ord_\ell(2)\mid g_{m+1}.
\end{equation}
\end{theorem}

\begin{proof}
Separate the first $m$ terms of \eqref{eq:Cn-Nn} after raising them to denominator
$2^{p_n}$, and apply the Euclidean algorithm.  For $n=m+2$,
$R_{m,m+2}=\pm(2^{g_{m+1}}-1)$, proving \eqref{eq:order-resonance}.
\end{proof}

\begin{corollary}[Fixed-lag logarithmic coprimality]\label{cor:fixed-lag-gcd}
\status{A+E}
For every fixed $r\ge1$,
\begin{equation}\label{eq:fixed-lag-gcd}
 \gcd(N_n,N_{n-r})
 \le |R_{n-r,n}|
 \le r\,2^{p_n-p_{n-r+1}}
 =N_n^{o(1)}.
\end{equation}
\end{corollary}

\begin{proof}
The first two inequalities follow from \eqref{eq:gcd-resonance} and
\eqref{eq:Rmn}.  The prime number theorem implies
$p_n-p_{n-r+1}=o(p_n)$ for fixed $r$.  Moreover
$N_n/2^{p_n}=C_n\to C_{PP}>0$, so
$\log N_n\sim p_n\log2$.
\end{proof}

\subsection{The third support series}

Because $N_1=N_2=1$, write
\begin{equation}\label{eq:FN}
 F_N(z):=\sum_{n\ge1}z^{N_n}=2z+\sum_{n\ge3}z^{N_n}.
\end{equation}

\begin{theorem}[Dyadic-numerator support]\label{thm:FN-values}
\status{A+L}
The unit circle is a natural boundary for $F_N$; in particular $F_N$ is transcendental
over $\mathbb C(z)$ and not D-finite.  For every algebraic $\alpha$ with
$0<|\alpha|<1$, the value $F_N(\alpha)$ is transcendental.  For every rational $r$
with $0<|r|<1$, the value $F_N(r)$ is a Liouville number.
\end{theorem}

\begin{proof}
For $n\ge4$,
\[
 \frac{N_n}{N_{n-1}}
 =2^{g_{n-1}}+\frac{(-1)^{n+1}}{N_{n-1}}
 \ge4-\frac15>1.
\]
After the duplicate initial exponent is removed, the Hadamard gap theorem gives the
natural boundary.  Corollary~5 of Corvaja--Zannier \cite{CorvajaZannier} gives
transcendence at nonzero algebraic points of the disk.  Finally the ratios have infinite
limsup because the prime gaps are unbounded, and the exponent differences tend to
infinity.  Apply \cref{lem:Liouville-lacunary}.
\end{proof}

\begin{corollary}[Arithmetic Picard fibres of $F_N$]
\label{cor:FN-arithmetic-Picard}
\status{L+E}
Let $a\in\overline{\mathbb Q}$, let $\zeta\in\partial\mathbb D$, and let $U$ be any
neighborhood of $\zeta$.  Then
\[
 \{z\in U\cap\mathbb D:F_N(z)=a\}
\]
contains infinitely many transcendental numbers.  In particular, every point of the unit
circle is an accumulation point of transcendental zeroes of $F_N$.
\end{corollary}

\begin{proof}
After combining the duplicated initial exponent, $F_N$ is an unbounded
Hadamard-lacunary series.  Murai's sectorial Picard theorem
\cite[Remark~21]{Murai} says that it assumes every complex value infinitely often in
every boundary sector.  The preimage of a value is discrete, so the infinitely many
solutions in a fixed sector leave every compact subset of $\mathbb D$.  Choosing first
an angular sector and then an outer annulus contained in $U$ proves the claim.  Any nonzero
algebraic solution would contradict \cref{thm:FN-values}; when $a=0$, the algebraic
solution $z=0$ accounts for at most one point.
\end{proof}

Combining \cref{thm:FQ-FH-algebraic,cor:FQ-FH-Liouville,thm:FN-values} gives the
uniform value statement
\begin{equation}\label{eq:three-series-values}
 \begin{array}{c|c|c}
 &0<|\alpha|<1,\ \alpha\in\overline{\mathbb Q}
 &0<|r|<1,\ r\in\mathbb Q\\ \hline
 F_Q,F_H,F_N&\text{transcendental}&\text{Liouville}.
 \end{array}
\end{equation}
These are values of support series, not the continued-fraction numbers in
\eqref{eq:mu-two}; the two irrationality-exponent statements concern different numbers.

\begin{theorem}[Hypertranscendence of the three support series]
\label{thm:three-series-hypertranscendental}
\status{L+E}
Each of $F_Q,F_H,F_N$ is hypertranscendental over $\mathbb C(z)$: none satisfies a
nonzero algebraic differential equation
\[
 P(z,F,F',\ldots,F^{(d)})=0,
 \qquad P\in\mathbb C[z,X_0,\ldots,X_d].
\]
\end{theorem}

\begin{proof}
The gap theorem of Lipshitz and Rubel \cite[Theorem~4.1]{LipshitzRubel} excludes
differential algebraicity for a nonpolynomial Hadamard-gap series.  It applies directly
to $F_Q$, because $Q_{m+1}>5Q_m$, and to $F_N$, whose successive distinct exponents
have ratio greater than $3$ from some point onward.

For $F_H$, let $\omega^3=1$ and use the root-of-unity filters
\[
 F_{H,r}(z):=\frac13\sum_{j=0}^2\omega^{-rj}F_H(\omega^jz),
 \qquad r=1,2.
\]
By \eqref{eq:H-mod3}, $F_H=F_{H,1}+F_{H,2}$.  Within either infinite filtered
support, successive exponents have ratio greater than $2$: if their original indices are
adjacent and the residues agree, the corresponding $h$ cannot equal $1$, while skipping
an index gives $H_{n+2}>2H_n$.  At least one filter is infinite, so it is a
Hadamard-gap series.  Differentially algebraic functions are closed under root-of-unity
substitution and finite sums.  If $F_H$ were differentially algebraic, both filters would
be, contradicting the gap theorem for the infinite one.
\end{proof}

Finally, \eqref{eq:Cn-Nn} proves the OEIS limit
\[
 \frac{N_n}{2^{p_n}}\longrightarrow C_{PP}.
\]
Together with \eqref{eq:2CPP-binary}, it also exposes the one-bit shift in the binary
expansion comment: A071986 is the digit word of $2C_{PP}$, whereas $C_{PP}$ has one
additional leading zero after the radix point.

\section{The Mellin boundary of the dyadic numerator series}
\label{sec:ZN-boundary}

Set $C=C_{PP}$.  By \eqref{eq:Cn-Nn},
\[
 N_n=C_n2^{p_n},\qquad C_n\longrightarrow C>0,
\]
and the alternating-series remainder gives
\begin{equation}\label{eq:Cn-tail}
 |C_n-C|\le2^{-p_{n+1}}.
\end{equation}
Attach to the numerator cocycle the Dirichlet series
\begin{equation}\label{eq:ZN-def}
 Z_N(s):=\sum_{n\ge1}N_n^{-s},\qquad \Re s>0,
\end{equation}
where positive real logarithms are used.  Since $\log N_n\sim p_n\log2$, its
abscissa of absolute convergence is $0$.

\subsection{Mellin representation and prime-support decomposition}

Let
\begin{equation}\label{eq:prime-support-series}
 P_{\rm pr}(z):=\sum_{p\in\PP}z^p,\qquad |z|<1.
\end{equation}

\begin{theorem}[Mellin representation]\label{thm:ZN-Mellin}
\status{E}
For $\Re s>0$,
\begin{equation}\label{eq:ZN-Mellin}
 \Gamma(s)Z_N(s)=\int_0^\infty t^{s-1}F_N(e^{-t})\,\dd t.
\end{equation}
\end{theorem}

\begin{proof}
For $\sigma=\Re s>0$, Tonelli's theorem applies because
\[
 \int_0^\infty t^{\sigma-1}\sum_{n\ge1}e^{-tN_n}\,\dd t
 =\Gamma(\sigma)\sum_{n\ge1}N_n^{-\sigma}<\infty.
\]
Termwise Mellin integration proves \eqref{eq:ZN-Mellin}; the multiplicity-two initial
exponent in $F_N$ agrees with $N_1=N_2=1$.
\end{proof}

\begin{proposition}[Prime-support decomposition]\label{prop:ZN-decomposition}
\status{E}
There is a function $\mathscr R$, holomorphic on $\Re s>-1$, such that
\begin{equation}\label{eq:ZN-decomposition}
 Z_N(s)=C^{-s}P_{\rm pr}(2^{-s})+\mathscr R(s)\qquad(\Re s>0),
\end{equation}
where
\begin{equation}\label{eq:ZN-R}
 \mathscr R(s)=\sum_{n\ge1}2^{-sp_n}(C_n^{-s}-C^{-s}).
\end{equation}
The latter series converges normally on compact subsets of $\Re s>-1$.
\end{proposition}

\begin{proof}
Let $K$ be compact in $\Re s>-1$ and put
$\sigma_0=\inf_{s\in K}\Re s>-1$.  The derivatives of $x\mapsto x^{-s}$ are
uniformly bounded for $s\in K$ on an interval containing $C$ and all late $C_n$.
Thus \eqref{eq:Cn-tail} gives
\[
 |2^{-sp_n}(C_n^{-s}-C^{-s})|
 \ll_K2^{-\sigma_0p_n-p_{n+1}}
 \le2^{-(1+\sigma_0)p_n}.
\]
The majorant is summable.  Normal convergence proves the assertion.
\end{proof}

\begin{theorem}[Natural boundary of $Z_N$]\label{thm:ZN-natural-boundary}
\status{A+E}
The imaginary axis $\Re s=0$ is a natural boundary for $Z_N$: there is no holomorphic
continuation across any point of that line.  Consequently, $Z_N$ has no meromorphic
continuation through a nonempty open subarc of it.
\end{theorem}

\begin{proof}
The integer-coefficient series $P_{\rm pr}$ has radius one and is not rational.  Indeed,
its bounded coefficient sequence cannot be eventually periodic: if it had eventual
period $q$, then a sufficiently large prime $p$ would force the composite number
$p+pq=p(1+q)$ to be prime.  P\'olya--Carlson \cite{Polya,Carlson} makes the unit
circle a natural boundary for $P_{\rm pr}$.

If $Z_N$ continued through a point of $\Re s=0$, then
$P_{\rm pr}(2^{-s})=C^s(Z_N(s)-\mathscr R(s))$ would continue there.  The map
$s\mapsto2^{-s}$ has nonzero derivative and a local inverse, contradicting the natural
boundary of $P_{\rm pr}$.  A meromorphic continuation through an open arc would be
holomorphic near a non-pole point of that arc and is therefore excluded as well.
\end{proof}

\subsection{The inverse-logarithmic Gumbel fingerprint}

For $m\ge0$, define
\begin{equation}\label{eq:Gumbel-moments}
 A_m:=(-1)^m\Gamma^{(m)}(1)
 =\int_0^\infty e^{-u}(-\log u)^m\,\dd u.
\end{equation}
The first values are
\begin{equation}\label{eq:first-Gumbel-moments}
 A_0=1,\quad A_1=\gamma,\quad
 A_2=\gamma^2+\frac{\pi^2}{6},\quad
 A_3=\gamma^3+\frac{\gamma\pi^2}{2}+2\zeta(3).
\end{equation}

\begin{theorem}[All-orders boundary expansion]\label{thm:ZN-Gumbel}
\status{A+E}
Let $y=s\log2$ and $\Lambda=\log(1/y)$, with $s>0$.  For every fixed $M\ge0$,
\begin{equation}\label{eq:ZN-Gumbel-expansion}
 yZ_N(s)=\sum_{m=0}^{M}\frac{A_m}{\Lambda^{m+1}}
 +O_M(\Lambda^{-M-2})\qquad(s\downarrow0).
\end{equation}
Equivalently, in the Poincar\'e sense,
\begin{equation}\label{eq:ZN-Gumbel-Poincare}
 (s\log2)Z_N(s)\sim
 \sum_{m\ge0}\frac{(-1)^m\Gamma^{(m)}(1)}
 {\bigl(\log(1/(s\log2))\bigr)^{m+1}}.
\end{equation}
\end{theorem}

\begin{proof}
Write $S(y)=P_{\rm pr}(e^{-y})=\sum_pe^{-yp}$.  The classical zero-free-region
form of the prime number theorem gives
\begin{equation}\label{eq:PNT-Li-error}
 \pi(x)=\operatorname{Li}(x)+O\bigl(xe^{-c\sqrt{\log x}}\bigr).
\end{equation}
Stieltjes integration and one integration by parts give, for some $c_1>0$,
\begin{equation}\label{eq:prime-Laplace-reduction}
 yS(y)=\int_{2y}^\infty\frac{e^{-u}}{\Lambda+\log u}\,\dd u
       +O\bigl(e^{-c_1\sqrt\Lambda}\bigr).
\end{equation}
Split the integral at $e^{-\Lambda/2}$ and $e^{\Lambda/2}$.  The exterior pieces are
smaller than every negative power of $\Lambda$.  In the middle range the finite
geometric expansion is
\[
 \frac1{\Lambda+\log u}
 =\sum_{m=0}^{M}\frac{(-\log u)^m}{\Lambda^{m+1}}
 +O_M\left(\frac{|\log u|^{M+1}}{\Lambda^{M+2}}\right).
\]
Every logarithmic moment against $e^{-u}\dd u$ is finite, so
\begin{equation}\label{eq:prime-Gumbel-expansion}
 yP_{\rm pr}(e^{-y})
 =\sum_{m=0}^{M}\frac{A_m}{\Lambda^{m+1}}+O_M(\Lambda^{-M-2}).
\end{equation}
Finally, $C^{-s}=1+O(s)$ and $\mathscr R$ is bounded near zero.  The resulting
corrections to $yZ_N(s)$ are smaller than every negative power of $\Lambda$, proving
\eqref{eq:ZN-Gumbel-expansion}.
\end{proof}

\begin{corollary}[Exact boundary-moment recovery]
\label{cor:ZN-moment-recovery}
\status{E}
For every $m\ge0$,
\begin{equation}\label{eq:ZN-moment-recovery}
 A_m=\lim_{s\downarrow0}\Lambda^{m+1}
 \left((s\log2)Z_N(s)-\sum_{j=0}^{m-1}\frac{A_j}{\Lambda^{j+1}}\right),
 \qquad \Lambda=\log\frac1{s\log2}.
\end{equation}
In particular,
\begin{equation}\label{eq:ZN-gamma-recovery}
 \gamma=\lim_{s\downarrow0}\Lambda^2
 \left((s\log2)Z_N(s)-\frac1\Lambda\right),
\end{equation}
and \cref{thm:ZN-Mellin} gives the equivalent $F_N$-formula
\begin{equation}\label{eq:FN-gamma-recovery}
 \gamma=\lim_{s\downarrow0}\Lambda^2\left[
 \frac{s\log2}{\Gamma(s)}\int_0^\infty t^{s-1}F_N(e^{-t})\,\dd t
 -\frac1\Lambda\right].
\end{equation}
\end{corollary}

\begin{proof}
Subtract the first $m$ terms in \cref{thm:ZN-Gumbel} and multiply the remainder by
$\Lambda^{m+1}$.
\end{proof}

\begin{remark}[Gumbel interpretation]\label{rem:Gumbel-interpretation}
If $U$ is exponential with mean one and $G=-\log U$, then $G$ has the standard
Gumbel distribution and
\[
 A_m=\mathbb E(G^m),\qquad
 \mathbb E(e^{tG})=\Gamma(1-t)\quad(t<1).
\]
Its cumulants are
\begin{equation}\label{eq:Gumbel-cumulants}
 \kappa_1=\gamma,
 \qquad \kappa_m=(m-1)!\zeta(m)\quad(m\ge2),
\end{equation}
because
\begin{equation}\label{eq:log-Gamma-cumulants}
 \log\Gamma(1-t)=\gamma t+\sum_{m\ge2}\frac{\zeta(m)}{m}t^m.
\end{equation}
Thus $\zeta(2)=A_2-A_1^2$ and
$2\zeta(3)=A_3-3A_2A_1+2A_1^3$.  These are boundary-coefficient
identities; they do not decide the irrationality of $\gamma$ or the transcendence of an
individual odd zeta value.
\end{remark}

\subsection{Boundary moments and a regularized Gamma determinant}

Let $\mathcal H=\ell^2(\mathbb N)^{\oplus5}$ and define
\begin{equation}\label{eq:A-Gamma}
 \mathsf A_\Gamma:=
 \operatorname{diag}\left(
 \frac1n,\frac1n,\frac1n,-\frac1n,-\frac1n:n\ge1\right).
\end{equation}
This compact self-adjoint operator is Hilbert--Schmidt but not trace class.

\begin{theorem}[Boundary--Gamma determinant correspondence]
\label{thm:boundary-Gamma-determinant}
\status{E}
For the truncation to $n\le N$,
\begin{equation}\label{eq:A-Gamma-FP-trace}
 \FPTr\mathsf A_\Gamma:=
 \lim_{N\to\infty}(\Tr\mathsf A_{\Gamma,N}-\log N)=A_1=\gamma,
\end{equation}
whereas, for $m\ge2$,
\begin{equation}\label{eq:A-Gamma-traces}
 \Tr\mathsf A_\Gamma^m=(3+2(-1)^m)\zeta(m)
 =\frac{3+2(-1)^m}{(m-1)!}\,\kappa_m.
\end{equation}
Moreover,
\begin{equation}\label{eq:A-Gamma-det2}
 \det\nolimits_2(I-t\mathsf A_\Gamma)
 =\frac{e^{\gamma t}}{\Gamma(1-t)^3\Gamma(1+t)^2}.
\end{equation}
If $\mathcal B(t)=\sum_{m\ge0}A_mt^m/m!=\Gamma(1-t)$, then
\begin{equation}\label{eq:boundary-det-correspondence}
 \mathcal B(t)^3\mathcal B(-t)^2
 =e^{A_1t}\det\nolimits_2(I-t\mathsf A_\Gamma)^{-1}.
\end{equation}
Thus the inverse-logarithmic fingerprint of $Z_N$ and this regularized trace tower
determine one another.
\end{theorem}

\begin{proof}
$\Tr\mathsf A_{\Gamma,N}=(3-2)H_N$, which proves
\eqref{eq:A-Gamma-FP-trace}; absolute summation proves
\eqref{eq:A-Gamma-traces}.  The two Weierstrass products
\[
 \prod_{n\ge1}(1-t/n)e^{t/n}=\frac{e^{\gamma t}}{\Gamma(1-t)},
 \qquad
 \prod_{n\ge1}(1+t/n)e^{-t/n}=\frac{e^{-\gamma t}}{\Gamma(1+t)}
\]
give \eqref{eq:A-Gamma-det2}; see \cite{SimonTraceIdeals}.  The
moment--cumulant transformations are mutually inverse universal polynomial maps, which
proves the final assertion.
\end{proof}

\subsection{Mahler exclusion and algebraic separation}

\begin{theorem}[All-base non-Mahler property]\label{thm:FN-non-Mahler}
\status{L+E}
For no integer $k\ge2$ is $F_N(z)$ a $k$-Mahler function.
\end{theorem}

\begin{proof}
Write $F_N(z)=\sum_{m\ge0}a_mz^m$.  If its finite-valued coefficient sequence were
$k$-automatic, the canonical base-$k$ representations of the support would form an
infinite regular language.  The pumping lemma forces such a language to have at least
$cL$ words of length at most $L$.  Here, however,
\begin{equation}\label{eq:FN-support-count}
 \#\{n:N_n<k^L\}
 \sim\frac{\log k}{\log2}\frac{L}{\log L}=o(L),
\end{equation}
because $\log N_n\sim p_n\log2$.  Thus $(a_m)$ is not $k$-automatic.  A
finite-valued $k$-Mahler series in characteristic zero is $k$-automatic by
Adamczewski--Bell--Smertnig \cite[Theorem~11.1]{AdamczewskiBellSmertnig}, giving the
contradiction.
\end{proof}

For a rational $r$ with $0<|r|<1$, put $L_r:=F_N(r)$.

\begin{theorem}[Liouville--$\pi$ algebraic independence]
\label{thm:FN-pi-independence}
\status{L+E}
For every such $r$,
\begin{equation}\label{eq:FN-pi-independence}
 \operatorname{trdeg}_{\overline{\mathbb Q}}
 \overline{\mathbb Q}(L_r,\pi)=2.
\end{equation}
Consequently $L_r$ is algebraically independent from every $\zeta(2m)$, and
\[
 \overline{\mathbb Q}(L_r,\zeta(2),\zeta(4),\ldots)
 =\overline{\mathbb Q}(L_r,\pi^2)
\]
has transcendence degree $2$.
\end{theorem}

\begin{proof}
By \cref{thm:FN-values}, $L_r$ is a Liouville number and hence a Mahler
$U$-number.  Mahler's classification theorem says that algebraically dependent
transcendental numbers have the same class, while $\pi$ is an $S$- or $T$-number;
see \cite[Theorems~2.3 and~2.5]{ChalebgwaMorris}.  This proves
\eqref{eq:FN-pi-independence}.  Euler's formula expressing $\zeta(2m)$ as a nonzero
rational multiple of $\pi^{2m}$ gives the rest.
\end{proof}

For example, the explicit number
\begin{equation}\label{eq:explicit-FN-pi-number}
 \frac6{\pi^2}\left(
 1+2^{-5}+2^{-19}+2^{-305}+2^{-1219}+\cdots\right)
 =\frac{6F_N(1/2)}{\pi^2}
\end{equation}
is transcendental.

\begin{theorem}[Relative algebraic independence of boundary jets]
\label{thm:relative-Gamma-jets}
\status{L+E}
Fix $r\in\mathbb Q$, $0<|r|<1$, and put
\[
 \tau_r(N):=\#\left\{1\le m\le N:
 \operatorname{trdeg}_{\overline{\mathbb Q}}
 \overline{\mathbb Q}(L_r,A_m)=2\right\}.
\]
Then
\begin{equation}\label{eq:relative-Gamma-lower-bound}
 \tau_r(N)\ge
 \left\lceil\sqrt{2\lfloor N/2\rfloor+\frac94}-\frac32\right\rceil.
\end{equation}
In particular, infinitely many Gumbel boundary coefficients are algebraically independent
from $L_r$.
\end{theorem}

\begin{proof}
Put $K=\overline{\mathbb Q}(L_r)$ and $M=\lfloor N/2\rfloor$.  The reflection
formula $\Gamma(1-t)\Gamma(1+t)=\pi t/\sin\pi t$ gives, for $0\le k\le M$,
\begin{equation}\label{eq:Gamma-quadratic-reflection}
 \sum_{j=0}^{2k}\frac{(-1)^jA_jA_{2k-j}}{j!(2k-j)!}
 =c_k\pi^{2k},\qquad c_k\in\mathbb Q^\times.
\end{equation}
Let $I$ consist of those indices $1\le j\le N$ for which $A_j$ is transcendental
over $K$, and write $t=|I|$.  Adjoining the remaining $A_j$ makes only a finite
algebraic extension $E/K$.  By \cref{thm:FN-pi-independence}, the numbers
$1,\pi^2,\ldots,\pi^{2M}$ are linearly independent over $E$.
Equations \eqref{eq:Gamma-quadratic-reflection} place them in the $E$-space of
polynomials of total degree at most two in the $t$ elements indexed by $I$, whose
dimension is at most $\binom{t+2}{2}$.  Hence
$M+1\le(t+1)(t+2)/2$, which rearranges to
\eqref{eq:relative-Gamma-lower-bound}.  This is the relative version of the standard
reflection-and-dimension argument for Gamma derivatives; compare \cite{PowersGamma}.
\end{proof}

\begin{corollary}[The first four boundary coefficients]
\label{cor:first-four-relative-Gamma}
\status{E}
At least two of $A_1,A_2,A_3,A_4$ are algebraically independent from $L_r$.
In particular, at least one of $A_2,A_3,A_4$ is algebraically independent from $L_r$.
\end{corollary}

\begin{proof}
Besides $A_2-A_1^2=\zeta(2)$, the first four moments satisfy
\[
 5A_4-20A_1A_3-27A_2^2+84A_1^2A_2-42A_1^4=0.
\]
This follows by substituting the complete Bell-polynomial expressions for the Gamma
derivatives and using $\zeta(4)=\tfrac25\zeta(2)^2$.  If exactly one of
$A_1,A_2,A_3,A_4$ were transcendental over $\overline{\mathbb Q}(L_r)$, the displayed
relation is a nonconstant polynomial in that element (for $A_3$, use
$A_1=\gamma\ne0$), making it algebraic as well; if none were transcendental, all four
would already be algebraic.  Then $A_2-A_1^2=\pi^2/6$ would make $\pi$ algebraic over
$\overline{\mathbb Q}(L_r)$, contradicting \cref{thm:FN-pi-independence}.
\end{proof}

\section{The alternating prime-index Dirichlet series}

Define
\begin{equation}\label{eq:Palt}
 \Palt(s):=\sum_{k\ge1}(-1)^{k+1}p_k^{-s}
 =\sum_{j\ge1}\bigl(p_{2j-1}^{-s}-p_{2j}^{-s}\bigr),
\end{equation}
and let $\xi(x)$ be the indicator of
$\bigcup_{j\ge1}[p_{2j-1},p_{2j})$.

\subsection{Laplace--Mellin representation}

\begin{theorem}[Analyticity and complete monotonicity]\label{thm:Palt-integral}
\status{E}
The function $\Palt$ is holomorphic on $\Re s>0$ and
\begin{equation}\label{eq:Palt-integral}
 \Palt(s)=s\int_2^\infty\xi(x)x^{-s-1}\,\dd x.
\end{equation}
For real $\sigma>0$ it is positive, and
$F(\sigma)=\Palt(\sigma)/\sigma$ is strictly completely monotone.
\end{theorem}

\begin{proof}
For each pair,
\[
 p_{2j-1}^{-s}-p_{2j}^{-s}
 =s\int_{p_{2j-1}}^{p_{2j}}x^{-s-1}\,\dd x.
\]
The disjoint intervals make the sum locally uniformly dominated on every half-plane
$\Re s\ge\delta>0$, proving \eqref{eq:Palt-integral}.  With $x=e^t$,
\[
 (-1)^rF^{(r)}(\sigma)
 =\int_{\log2}^{\infty}\xi(e^t)t^re^{-\sigma t}\,\dd t>0.
\]
\end{proof}

\subsection{A critical rank-one inverse spectrum}

Put
\[
 \alpha_j=p_{2j-1}^{-1},\qquad \beta_j=p_{2j}^{-1},
\]
so that $\alpha_1>\beta_1>\alpha_2>\beta_2>\cdots\downarrow0$, and define
\begin{equation}\label{eq:prime-perturbation-product}
 \Psi(z):=\prod_{j\ge1}\frac{1-z\beta_j}{1-z\alpha_j}.
\end{equation}
The product converges locally uniformly away from its poles because
$\sum_j(\alpha_j-\beta_j)<\infty$.

\begin{theorem}[Critical rank-one realization of prime interlacing]
\label{thm:prime-rank-one}
\status{A+E}
There are positive weights
\begin{equation}\label{eq:prime-rank-one-weights}
 w_j=(\alpha_j-\beta_j)
 \prod_{k\ne j}\frac{1-\beta_k/\alpha_j}{1-\alpha_k/\alpha_j}
\end{equation}
such that
\begin{equation}\label{eq:prime-pick-expansion}
 \Psi(z)=1+z\sum_{j\ge1}\frac{w_j}{1-z\alpha_j},
\end{equation}
and
\begin{equation}\label{eq:critical-weight-sums}
 \sum_jw_j=\Palt(1),\qquad
 \sum_j\frac{w_j}{\alpha_j}=1,
 \qquad
 \sum_j\frac{w_j}{\alpha_j^2}=\infty.
\end{equation}
Consequently $\Psi$ is a meromorphic Pick function.

On $\ell^2(\mathbb N)$ let
\[
 \mathsf A_{\rm o}=\operatorname{diag}(\alpha_1,\alpha_2,\ldots),\qquad
 v=(\sqrt{w_j})_{j\ge1},\qquad
 \mathsf A_{\rm e}=\mathsf A_{\rm o}-v\otimes v.
\]
Then $\mathsf A_{\rm o}$ and $\mathsf A_{\rm e}$ are positive compact injective
operators and
\begin{equation}\label{eq:prime-rank-one-spectra}
 \sigma(\mathsf A_{\rm o})=\{0\}\cup\{p_{2j-1}^{-1}:j\ge1\},\qquad
 \sigma(\mathsf A_{\rm e})=\{0\}\cup\{p_{2j}^{-1}:j\ge1\}.
\end{equation}
For $z$ away from the poles,
\begin{equation}\label{eq:prime-relative-determinant}
 \Psi(z)=\det\bigl((I-z\mathsf A_{\rm e})(I-z\mathsf A_{\rm o})^{-1}\bigr).
\end{equation}
Although neither operator is trace class, for every $\Re s>0$ one has
$\mathsf A_{\rm o}^s-\mathsf A_{\rm e}^s\in\mathcal S_1$ and
\begin{equation}\label{eq:prime-relative-trace}
 \Tr(\mathsf A_{\rm o}^s-\mathsf A_{\rm e}^s)
 =\sum_{j\ge1}(\alpha_j^s-\beta_j^s)=\Palt(s),
\end{equation}
where $x^s=e^{s\log x}$ for $x>0$, $0^s=0$, and the series is paired as displayed.

Finally, if $\mathsf A_{{\rm e},\tau}=\mathsf A_{\rm o}-\tau v\otimes v$, then
\begin{equation}\label{eq:critical-coupling}
 \mathsf A_{{\rm e},\tau}\ge0\quad(0\le\tau\le1),
 \qquad
 \mathsf A_{{\rm e},\tau}\text{ has exactly one negative eigenvalue}\quad(\tau>1).
\end{equation}
At $\tau=1$, zero is a continuous-spectrum threshold, not an eigenvalue.
\end{theorem}

\begin{proof}
For a finite interlacing product, partial fractions give
$\Psi_J(z)=1+z\sum_{j\le J}w_{j,J}/(1-z\alpha_j)$ with positive residues.
Passing monotonically to the limit gives \eqref{eq:prime-rank-one-weights} and
\eqref{eq:prime-pick-expansion}.  The finite identities
\[
 \sum_{j\le J}w_{j,J}=\sum_{j\le J}(\alpha_j-\beta_j),
 \qquad
 1-\sum_{j\le J}\frac{w_{j,J}}{\alpha_j}
 =\prod_{j\le J}\frac{\beta_j}{\alpha_j}
\]
give the first two sums in \eqref{eq:critical-weight-sums}; the last product tends to
zero because the odd and even prime blocks both have positive lower density
\cite{ChungLi}.

For $t>0$ write
\[
 D_o(t)=\sum_j\int_{p_{2j-1}}^{p_{2j}}\frac{t\,\dd u}{u(u+t)}=-\log\Psi(-t)
\]
and define $D_e(t)$ using the complementary intervals
$[p_{2j},p_{2j+1}]$.  Then
$D_o(t)+D_e(t)=\log(1+t/2)$, while both terms tend to infinity by the same
positive-density theorem.  Hence
$t\Psi(-t)=2t(t+2)^{-1}e^{D_e(t)}\to\infty$.  Since
\[
 \Psi(-t)=\sum_j\frac{w_j/\alpha_j}{1+t\alpha_j},
\]
monotone convergence proves the last sum in
\eqref{eq:critical-weight-sums}.

Cauchy--Schwarz now gives
\[
 |\langle u,v\rangle|^2\le
 \langle \mathsf A_{\rm o}u,u\rangle\sum_j\frac{w_j}{\alpha_j}
 =\langle \mathsf A_{\rm o}u,u\rangle,
\]
so $\mathsf A_{\rm e}\ge0$.  Equality would require $u$ proportional to
$\mathsf A_{\rm o}^{-1}v$, but
$\|\mathsf A_{\rm o}^{-1}v\|^2=\infty$, proving injectivity.  The matrix determinant
lemma gives
\eqref{eq:prime-relative-determinant}; its interlacing zeros give
\eqref{eq:prime-rank-one-spectra}.  The same quadratic-form estimate proves positivity
for $\tau\le1$, while finite truncations of $\mathsf A_{\rm o}^{-1}v$ give one negative direction when
$\tau>1$; a rank-one subtraction gives at most one.

The Krein spectral-shift function is
\[
 \eta(x)=\one_{\cup_j(\beta_j,\alpha_j)}(x)=\xi(1/x)\quad\text{a.e.}
\]
For $u>0$ put
\[
 R_u=(\mathsf A_{\rm e}+uI)^{-1}-(\mathsf A_{\rm o}+uI)^{-1},
 \qquad d(u)=\Psi(-1/u).
\]
Sherman--Morrison and \eqref{eq:prime-pick-expansion} give
\begin{equation}\label{eq:prime-resolvent-rank-one}
 R_u=\frac{(\mathsf A_{\rm o}+uI)^{-1}v\otimes
                 (\mathsf A_{\rm o}+uI)^{-1}v}{d(u)},
 \qquad
 \lVert R_u\rVert_1=\frac{d'(u)}{d(u)}.
\end{equation}
The estimate $T\Psi(-T)\to\infty$ proved above gives $d(u)/u\to\infty$ as
$u\downarrow0$, and hence $|\log d(u)|=O(|\log u|)$; also
$\lVert R_u\rVert_1=O(u^{-2})$ as $u\to\infty$.  Integration by parts therefore shows
\[
 \int_0^\infty u^\sigma\lVert R_u\rVert_1\,\dd u<\infty
 \qquad(0<\sigma<1).
\]
For $0<\Re s<1$, the trace-norm convergent fractional-power formula
\begin{equation}\label{eq:prime-fractional-resolvent}
 \mathsf A_{\rm o}^s-\mathsf A_{\rm e}^s
 =\frac{\sin\pi s}{\pi}\int_0^\infty u^sR_u\,\dd u
\end{equation}
proves membership in $\mathcal S_1$.  Differentiating
$d(u)=\prod_j(u+\beta_j)/(u+\alpha_j)$ and integrating termwise gives
\eqref{eq:prime-relative-trace} in this strip.  The identity
\[
 \mathsf A_{\rm o}^{2s}-\mathsf A_{\rm e}^{2s}
 =(\mathsf A_{\rm o}^s-\mathsf A_{\rm e}^s)\mathsf A_{\rm o}^s
 +\mathsf A_{\rm e}^s(\mathsf A_{\rm o}^s-\mathsf A_{\rm e}^s)
\]
propagates trace-class membership to every $\Re s>0$, and analytic continuation
propagates the trace identity.  This is the concrete resolvent form of the usual
spectral-shift argument; compare
\cite{BirmanYafaevSSF,IonascuRankOne,LiawRankOne}.
\end{proof}

\begin{corollary}[Prime subordinator and a strict moment inequality]
\label{cor:prime-subordinator}
\status{A+E}
$\Psi(-t)$ is the Laplace transform of an infinitely divisible random variable
$X_{\PP}\ge0$ with L\'evy measure
\begin{equation}\label{eq:prime-Levy-measure}
 \frac1u\sum_{j\ge1}
 (e^{-p_{2j-1}u}-e^{-p_{2j}u})\,\dd u.
\end{equation}
It has infinite activity, finite mean, and
\[
 \kappa_m(X_{\PP})=(m-1)!\Palt(m)\qquad(m\ge1).
\]
In particular,
\begin{equation}\label{eq:Palt-Hankel-inequality}
 4\Palt(1)\Palt(3)-3\Palt(2)^2-\Palt(1)^4>0.
\end{equation}
\end{corollary}

\begin{proof}
Frullani's formula gives the L\'evy--Khintchine representation of
$\log\Psi(-t)$.  The activity is infinite because
$\sum_j\log(p_{2j}/p_{2j-1})=\infty$, and the mean is $\Palt(1)$.
Differentiation at zero gives the cumulants.  Finally, strict positivity of the
$2\times2$ Hankel determinant for the infinite-support measure
$\sum_jw_j\delta_{\alpha_j}$ reduces algebraically to
\eqref{eq:Palt-Hankel-inequality}; compare \cite{SchillingBernstein}.
\end{proof}

\subsection{The exact right edge in the half-plane of absolute convergence}

Let $\cP(\sigma)=\sum_pp^{-\sigma}$ be the prime zeta function for $\sigma>1$.
Define $\sigma_{\mathrm e}$ by either of the equivalent equations
\begin{equation}\label{eq:sigma-edge}
 \sum_{p\ge3}\left(\frac2p\right)^{\sigma_{\mathrm e}}=1,
 \qquad
 \cP(\sigma_{\mathrm e})=2^{1-\sigma_{\mathrm e}}.
\end{equation}
The solution is unique, and numerically
\begin{equation}\label{eq:sigma-edge-numeric}
 \sigma_{\mathrm e}
 =1.7795446535469941164458987869655\ldots.
\end{equation}

Let
\[
 Z_{>1}:=\{\Re\rho:\Palt(\rho)=0,\ \Re\rho>1\}.
\]

\begin{theorem}[Exact zero-projection edge]\label{thm:Palt-edge}
\status{L+E}
The closure of $Z_{>1}$ relative to $(1,\infty)$ is
\begin{equation}\label{eq:Palt-relative-closure}
 \overline{Z_{>1}}^{\,(1,\infty)}=(1,\sigma_{\mathrm e}],
\end{equation}
or, equivalently, its ordinary closure in $\mathbb R$ is
\begin{equation}\label{eq:Palt-ordinary-closure}
 \overline{Z_{>1}}^{\,\mathbb R}=[1,\sigma_{\mathrm e}].
\end{equation}
Moreover
\begin{equation}\label{eq:Palt-zero-free-edge}
 \Palt(s)\ne0\qquad(\Re s\ge\sigma_{\mathrm e}).
\end{equation}
Thus $\sup Z_{>1}=\sigma_{\mathrm e}$, but the supremum is not attained.
\end{theorem}

\begin{proof}
On $\Re s>1$, \eqref{eq:Palt} is absolutely convergent and uniformly almost periodic
on every closed interior strip.  Its frequencies $-\log p$ are linearly independent over
$\mathbb Q$: after clearing denominators, any rational relation would contradict unique
factorization.

The polygon criterion of Sepulcre--Vidal \cite[Theorem~6]{SepulcreVidal} says that
$\sigma>1$ belongs to the closure of zero real projections exactly when every side
$p^{-\sigma}$ is at most the sum of the remaining sides.  Since $2^{-\sigma}$ is the
largest, the only nonautomatic inequality is
\[
 2^{-\sigma}\le\sum_{p\ge3}p^{-\sigma},
\]
which holds exactly for $1<\sigma\le\sigma_{\mathrm e}$.  This proves the two closure
statements.

For $\sigma>\sigma_{\mathrm e}$, the tail has modulus strictly smaller than
$2^{-\sigma}$, so no zero is possible.  Suppose a zero existed at
$s=\sigma_{\mathrm e}+it$.  Equality would then hold in the triangle inequality, forcing
all tail vectors to align.  Comparing the terms for $3,5$ and for $5,7$ gives
\[
 t\log(5/3)\equiv\pi\pmod{2\pi},
 \qquad
 t\log(7/5)\equiv\pi\pmod{2\pi}.
\]
Their logarithm ratio would be rational, which would force an impossible multiplicative
relation between $5/3$ and $7/5$.  Hence the boundary is also zero-free.
\end{proof}

\begin{theorem}[Linear order of zeros in the absolute-convergence strip]
\label{thm:Palt-AP-zero-count}
\status{L+E}
For every $1<a<b<\sigma_{\mathrm e}$, counting multiplicity,
\begin{equation}\label{eq:Palt-AP-zero-count}
 \#\{\rho:\Palt(\rho)=0,\ a<\Re\rho<b,\ 0<\Im\rho<T\}
 \asymp_{a,b}T.
\end{equation}
\end{theorem}

\begin{proof}
The closure theorem supplies an actual zero in every such open strip.  Choose a small
closed disk around one zero, with boundary free of zeros.  Relatively dense almost
periods approximate $\Palt$ uniformly on the boundary; Rouch\'e's theorem therefore
reproduces a zero in linearly many disjoint vertical translates, giving the lower bound.

For the upper bound, the compact vertical-translation hull consists of limits
\[
 g(s)=\sum_{k\ge1}(-1)^{k+1}\omega_kp_k^{-s},\qquad |\omega_k|=1.
\]
At any fixed $B>\sigma_{\mathrm e}$, first-term domination makes every hull element
nonzero at $B$, so none is identically zero.  Compactness, Hurwitz's theorem, and a
slightly enlarged unit-height rectangle give a uniform bound for the number of zeros in
each unit box.  Summing the boxes proves the upper bound.
\end{proof}

The theorem determines linear order, not an asymptotic constant, and it does not say that
every individual vertical line contains a zero.

\subsection{Strong universality below the line of absolute convergence}

Bayart and Kouroupis prove that
$\sum_{n\ge1}(-1)^np_n^{-s}=-\Palt(s)$ is strongly universal on
$1/2<\Re s<1$ \cite[Corollary 1.4]{BayartKouroupis}.  The following consequence is therefore a
literature application, not a new universality theorem.

\begin{corollary}[Universality zeros]\label{cor:Palt-universal-zeros}
\status{L+E}
For every $1/2<a<b<1$,
\begin{equation}\label{eq:Palt-universal-zeros}
 \#\{\rho:\Palt(\rho)=0,\ a<\Re\rho<b,\ 0<\Im\rho<T\}
 \gg_{a,b}T.
\end{equation}
\end{corollary}

\begin{proof}
Choose a closed disk $K$ in the strip, centered at $s_0$, and use the target
$f(s)=s-s_0$.  Strong universality gives a positive lower density of vertical shifts on
which $\Palt(s+i\tau)$ approximates $f$ more closely on $\partial K$ than
$|f|$.  Rouch\'e's theorem puts a zero in every translated disk.  A fixed zero can
correspond to a $\tau$-interval of length at most the disk diameter, so positive measure
of good shifts requires linearly many distinct zeros.
\end{proof}

\subsection{The Jessen zero-measure limit}

\begin{theorem}[Jessen zero measure for $\Palt$]
\label{thm:Palt-Jessen-measure}
\status{L+E}
For every $\sigma>1/2$, the limit
\begin{equation}\label{eq:Palt-Jessen-function}
 J(\sigma):=\lim_{T\to\infty}\frac1T
 \int_0^T\log|\Palt(\sigma+it)|\,\dd t
\end{equation}
exists.  The function $J$ is finite, continuous, and convex on $(1/2,\infty)$.
Writing a zero as $\rho=\beta+i\gamma$ and counting its multiplicity $m_\rho$, set
\[
 \nu_T:=\frac1T
 \sum_{\substack{\Palt(\rho)=0\\0<\gamma<T,\ \beta>1/2}}
 m_\rho\delta_\beta.
\]
Then, as locally finite measures on $(1/2,\infty)$,
\begin{equation}\label{eq:Palt-Jessen-vague}
 \nu_T\xrightarrow{\rm v}\frac1{2\pi}J'',
\end{equation}
where $J''$ is the distributional second derivative.  If $J$ is differentiable at
$a$ and $b$, with $1/2<a<b$, then
\[
 N(a,b;T):=
 \sum_{\substack{\Palt(\rho)=0\\a<\Re\rho<b,\ 0<\Im\rho<T}}m_\rho
\]
satisfies
\begin{equation}\label{eq:Palt-exact-zero-density}
 N(a,b;T)=\frac{J'(b)-J'(a)}{2\pi}T+o(T).
\end{equation}
The support of $J''$ in this half-line is exactly $(1/2,\sigma_{\rm e}]$.
\end{theorem}

\begin{proof}
Let $P_N(s)=\sum_{k\le N}(-1)^{k+1}p_k^{-s}$.  These Dirichlet polynomials are
analytic Bohr almost-periodic functions.  For every $1/2<a<b$, Landau's mean-square
theorem \cite[Part~24, Satz~37]{Landau1909} applied to the tails gives
\begin{align*}
 \lim_{T\to\infty}\frac1T\int_0^T\int_a^b
 |\Palt(\sigma+it)-P_N(\sigma+it)|^2\,\dd\sigma\,\dd t
 =\sum_{k>N}\int_a^bp_k^{-2\sigma}\,\dd\sigma,
\end{align*}
and the right side tends to zero.  This is precisely the area-mean approximation
hypothesis in the generalized Jessen theorem of Borchsenius and Jessen
\cite[Theorem~1]{BorchseniusJessen1948}; see also the ordinary almost-periodic theorem
of Jessen and Tornehave \cite{JessenTornehave1945}.  Moreover, on every closed strip
$1<\alpha_0\le\Re s\le\beta_0<\infty$, absolute convergence gives
$P_N\to\Palt$ uniformly.  Thus the Borchsenius--Jessen theorem applies with mean index
$p=2$ and yields existence,
continuity and convexity of $J$ and the interval formula
\eqref{eq:Palt-exact-zero-density} at differentiability endpoints.  Since these
endpoints are dense, the interval formula is equivalent to the vague convergence
\eqref{eq:Palt-Jessen-vague}.

Given a nonempty open interval $I\Subset(1/2,1)$, choose differentiability points
$a<b$ of $J$ in $I$.  Then \cref{cor:Palt-universal-zeros} and
\eqref{eq:Palt-exact-zero-density} give $J''((a,b))>0$; the same argument with
\cref{thm:Palt-AP-zero-count} applies when $I\Subset(1,\sigma_{\rm e})$.  Thus every
such interval meets the support.  The zero-free region
\eqref{eq:Palt-zero-free-edge} gives $J''((\sigma_{\rm e},\infty))=0$, intervals
approaching $\sigma_{\rm e}$ from the left have positive mass, and every neighborhood
of $1$ meets one of the preceding regions.  Hence the support relative to
$(1/2,\infty)$ is exactly $(1/2,\sigma_{\rm e}]$.
\end{proof}

\begin{remark}
Below $\sigma=1$ this is not a direct application of the ordinary
Jessen--Tornehave Dirichlet-series theorem.  Kronecker approximation gives
$\sigma_u(\Palt)=\sigma_a(\Palt)=1$; the extension to $\sigma>1/2$ specifically
uses the Borchsenius--Jessen area-mean theorem.  Also, $J''$ is a Radon measure, not
necessarily an ordinary density function.
\end{remark}

\subsection{Mellin--Barnes recovery of the parity prime constant}

\begin{theorem}\label{thm:Mellin-Barnes-CPP}
\status{E}
For every $c>0$,
\begin{equation}\label{eq:Mellin-Barnes-CPP}
 2C_{PP}=\frac1{\pi i}
 \int_{c-i\infty}^{c+i\infty}
 \Gamma(s)(\log2)^{-s}\Palt(s)\,\dd s.
\end{equation}
\end{theorem}

\begin{proof}
For $c>1$, insert $x=p_k\log2$ into Mellin inversion for $e^{-x}$ and interchange the
absolutely convergent sum and integral.  This gives the same formula with $C_{PP}$ and
factor $1/(2\pi i)$.  On closed strips in $\Re s>0$, \eqref{eq:Palt-integral} gives
$|\Palt(s)|\ll(1+|\Im s|)$, whereas $\Gamma(s)$ decays exponentially.  Contour shifting
is therefore valid to every $c>0$.  Multiplying by $2$ gives
\eqref{eq:Mellin-Barnes-CPP}.
\end{proof}

\section{The global parity frontier and Erd\H{o}s's alternating prime series}

The fixed-lag and multiplicatively weighted results above should not be confused with
unweighted cancellation of the prime-parity blocks.  Put
\begin{equation}\label{eq:Tpi}
 D_r(x):=\#\{n\le x:\pi(n)\equiv r\pmod2\},
 \qquad
 \mathcal T(x):=\sum_{n\le x}(-1)^{\pi(n)}=D_0(x)-D_1(x).
\end{equation}

Chung and Li proved the general unconditional estimate \cite{ChungLi}
\begin{equation}\label{eq:Chung-Li}
 \liminf_{x\to\infty}\frac1x
 \#\{n\le x:\pi(n)\equiv r\pmod t\}
 \ge\frac1{16t^2}.
\end{equation}
At $t=2$ this yields
\begin{equation}\label{eq:parity-positive-density}
 \liminf\frac{D_0(x)}x,\quad \liminf\frac{D_1(x)}x\ge\frac1{64},
 \qquad
 \limsup\frac{|\mathcal T(x)|}{x}\le\frac{31}{32}.
\end{equation}
Alboiu independently proved positive proportion in every residue class
\cite{Alboiu}.  Thus a simple Selberg argument giving, for example, $1/128$ would be
valid but weaker than known work.

\begin{proposition}[Thick excursions forced by large prime gaps]
\label{prop:thick-parity-excursions}
\status{L+E}
Let
\[
 G(X):=\max_{p_{k+1}\le X}(p_{k+1}-p_k).
\]
Then
\begin{equation}\label{eq:thick-parity-excursions}
 \max_{n\le X}|\mathcal T(n)|\ge\frac12G(X),
 \qquad
 \#\{n\le X:|\mathcal T(n)|>G(X)/4\}
 \ge\frac12G(X)-O(1).
\end{equation}
Consequently, for all sufficiently large $X$,
\begin{equation}\label{eq:FGKMT-parity-excursion}
 \max_{n\le X}|\mathcal T(n)|
 \gg\frac{\log X\,\log_2X\,\log_4X}{\log_3X}.
\end{equation}
\end{proposition}

\begin{proof}
On $[p_k,p_{k+1}-1]$ the sign is constant and
\[
 \mathcal T(p_{k+1}-1)-\mathcal T(p_k-1)=(-1)^kg_k.
\]
One endpoint therefore has modulus at least $g_k/2$.  Among the $g_k$ successive
integer values traversed on the block, at most $g_k/2+O(1)$ lie in
$[-g_k/4,g_k/4]$.  Maximize over the gaps below $X$ and apply \cite{FGKMT}.
\end{proof}

The parity equidistribution problem
\begin{equation}\label{eq:global-parity-open}
 \mathcal T(x)=o(x)
\end{equation}
remains open.  It is equivalent to
$D_0(x),D_1(x)\sim x/2$.  Notice that \eqref{eq:prime-parity-mu-lambda} proves
cancellation only after inserting $\mu$ or $\lambda$; it does not remove those weights.

\begin{theorem}[Punctured Fourier boundary]\label{thm:punctured-Erdos}
\status{A+E}
Let $\alpha>0$.  For every $|z|=1$ with $z\ne1$, the series
\begin{equation}\label{eq:punctured-series}
 \sum_{n\ge2}\frac{(-1)^{\pi(n)}z^n}{n(\log n)^\alpha}
\end{equation}
converges, and
\begin{equation}\label{eq:punctured-tail}
 \left|\sum_{n>M}\frac{(-1)^{\pi(n)}z^n}{n(\log n)^\alpha}\right|
 \ll_\alpha\frac1{|1-z|(\log M)^\alpha}.
\end{equation}
For $\alpha=1$, every boundary frequency except $z=1$ is therefore settled
unconditionally.
\end{theorem}

\begin{proof}
Summation by parts first gives the exact identity
\begin{equation}\label{eq:punctured-first-difference}
 (1-z)\sum_{n=1}^N\sigma_nz^n
 =z+2\sum_{p_k\le N}(-1)^kz^{p_k}-\sigma_Nz^{N+1}.
\end{equation}
The prime number theorem, even in its elementary upper-bound form, implies
\[
 \sum_{n\le N}\sigma_nz^n\ll\frac{N}{|1-z|\log N}.
\]
A second partial summation against $1/(n(\log n)^\alpha)$ proves
\eqref{eq:punctured-tail}.
\end{proof}

\begin{corollary}[The frequency $z=-1$]\label{cor:minus-one-parity}
\status{E}
For every $X\ge3$,
\begin{equation}\label{eq:minus-one-bounded-sums}
 \sum_{n\le X}(-1)^{n+\pi(n)}\in\{-3,-2,-1\}.
\end{equation}
Consequently $\sum_{n\ge1}(-1)^{n+\pi(n)}n^{-s}$ has abscissa of convergence $0$
and abscissa of absolute convergence $1$.
\end{corollary}

\begin{proof}
For $m\ge2$, the two terms at $2m-1$ and $2m$ have the same value of $\pi$ and
opposite ordinary parity, so every completed pair cancels.  The first two terms sum to
$-2$, and an unfinished pair changes this by $\pm1$.
\end{proof}

Tao \cite{TaoErdos} proved the exact equivalence between convergence of Erd\H{o}s's
series
\begin{equation}\label{eq:Erdos-series}
 \sum_{n\ge1}\frac{(-1)^n n}{p_n}
\end{equation}
and convergence of
\begin{equation}\label{eq:Tao-reduced-series}
 \sum_{m\ge2}\frac{(-1)^{\pi(m)}}{m\log m}.
\end{equation}
The qualitative estimate \eqref{eq:global-parity-open} would not by itself settle this
convergence: partial summation would leave a merely nonquantitative $o(1/(x\log x))$
integrand, which need not be integrable.  A sufficient estimate is
\begin{equation}\label{eq:E7-sufficient}
 \mathcal T(x)\ll
 \frac{x}{(\log\log x)^{1+\eps}}
\end{equation}
for some $\eps>0$; Tao displays the exponent $1.1$ and obtains the needed conclusion
under a strong Hardy--Littlewood prime-tuples hypothesis.

The honest hierarchy is therefore that published positive-proportion bounds are weaker
than $\mathcal T(x)=o(x)$, which in turn is weaker than integrable quantitative
cancellation such as \eqref{eq:E7-sufficient}.
The next useful input must detect a signed first moment or short-interval parity
correlation.  Merely bounding parity-separated nonnegative second moments of the gaps
does not force cancellation between odd- and even-indexed gap blocks.

\section{Limitations, open problems, and audited status}

\subsection{What the paper does not prove}

\begin{enumerate}
 \item \textbf{Twin primes and phase tomography.}  The curvature identity
 \eqref{eq:twin-curvature}, the modulo-$24$ sensor \eqref{eq:mod24-sensor}, the
 threshold \eqref{eq:single-lag-threshold}, and the phase decoder all re-encode finite
 prime-pair data.  The decoder recovers $G_{r,X}(h)$ from supplied phase correlations,
 but gives no lower bound or recurrence theorem and does not recover the temporal order
 of gaps from one fixed-$X$ correlation polynomial.  In particular, they do not prove
 $H_*=2$.
 \item \textbf{Gilbreath's conjecture.}  The affine inverse language, wheel criterion,
 corrected anchored right congruence, and transfer bounds describe obstruction words at
 fixed depth and height.  They do not prove that actual prime data avoid every obstruction.
 \item \textbf{Transcendence of $C_{PP}$.}  We prove irrationality and base-$2$
 non-normality, but neither the natural boundary of $B$ nor formal transcendence over
 $\Ftwo(z)$ decides the arithmetic nature of $B(1/2)$.
 \item \textbf{Algebraicity of the prime-gap continued fractions.}  Irrationality exponent
 $2$ is compatible with both algebraic and transcendental numbers.  The Liouville
 statements concern $F_Q(r),F_H(r),F_N(r)$, not $\alpha_P$ or $\beta$.
 \item \textbf{Euler's constant and odd zeta values.}  The Mellin boundary recovers
 $\gamma,\zeta(2),\zeta(3),\ldots$ as Gumbel moments, but boundary limits do not preserve
 arithmetic type.  No irrationality or transcendence result for $\gamma$, and no new
 individual transcendence result for an odd zeta value, follows.
 \item \textbf{Global parity cancellation.}  The M\"obius--Liouville orthogonality
 theorem, the Green inversion \eqref{eq:T-green}, and the punctured Fourier theorem do
 not imply \eqref{eq:global-parity-open}; even that open estimate would not by itself
 resolve Erd\H{o}s's series.
 \item \textbf{Certified computation.}  Numerical experiments and earlier ``all pass''
 reports without source and certificate bundles are not used as theorem inputs.
\end{enumerate}

\subsection{Concrete next problems}

\begin{enumerate}
 \item Close the gap
 \[
  \exp\!\left(\left(\frac1{16}-o(1)\right)(\log m)^2\right)\le p_b(m)
  \le\exp(Cm/\log m)
 \]
 unconditionally.  The CRT proof pinpoints the present bottleneck: Maynard's theorem
 supplies $\asymp\log k$ primes among $k$ candidates, whereas the conjectural optimum
 requires controlling a positive fraction or an arbitrary admissible subtuple.
 \item Improve the logarithmic rate in \cref{thm:quant-prime-block}, determine its sharp
 order, or obtain an effective theorem for general upper-Banach-sparse transition sets.
 The dynamical theorem \cref{thm:variation-orthogonality} remains qualitative outside
 the prime-block setting.
 \item Find genuinely quantitative input for the phase polynomials
 $\mathcal F_{X,h}(z)$.  Fourier inversion is lossless decoding, not a recurrence theorem,
 and fixed-$X$ data do not encode gap order.
 \item Determine whether the finite-state Gilbreath quotient can be reduced
 polynomially in $(H,r)$ while retaining actual-prime congruence data, or find an
 arithmetic invariant that separates surviving formal paths from genuine prime paths.
 \item Determine algebraicity or transcendence of $C_{PP}$ and of the three prime-gap
 continued fractions, and determine individual odd Gumbel cumulants.  New functional
 equations or value-transcendence input would be required; natural boundaries and exact
 boundary recovery alone are insufficient.
 \item In the Dirichlet problem, determine an explicit formula or finer regularity for the
 Jessen function $J$: in particular, decide whether $J''$ has atoms or an absolutely
 continuous component and describe zero statistics on individual vertical lines.
 \item For the Erd\H{o}s problem, first prove genuine cancellation
 $\mathcal T(x)=o(x)$, then strengthen it to an integrable rate such as
 \eqref{eq:E7-sufficient}.  By \eqref{eq:T-moment-equivalence}, the first step is exactly
 a signed first-moment asymptotic for the all-lag correlation; a viable analytic target is
 unconditional decay of a short-interval Fourier coefficient of prime-gap index parity.
\end{enumerate}

\subsection{Claim-status summary}

\begin{center}
\small
\begin{tabularx}{\textwidth}{@{}>{\raggedright\arraybackslash}p{0.27\textwidth}
 >{\raggedright\arraybackslash}p{0.49\textwidth}c@{}}
\toprule
Topic & Strongest statement used here & Status\\
\midrule
Factor complexity &
$\log p_b(m)\ge(1/16-o(1))(\log m)^2$ and
$\log p_b(m)=O(m/\log m)$; conditional optimal order & A/L/C\\
Orthogonality &
General qualitative sparse-transition theorem; uniform
$O(x\log_3x/\log_2x)$ prime-block rate & L/E\\
Local word laws &
Prime-rooted two-wall limit and reciprocal three-phase occurrence law & A/E\\
Finite diffraction &
Exact recovery of $\pi(N)$ and a renormalized high-pass Haar limit & A/E\\
Correlation decoding &
All-lag Green inversion, exact parity reconstruction, phase tomography, and $H_*$ recovery & E\\
Gauss--continuant geometry &
Common gap decoder, growth laws, irrationality exponent $2$, and $\dim_H Z\le2/3$ & A/E\\
Support series &
Natural boundaries, hypertranscendence, and arithmetic value classifications & L/E\\
Gilbreath language &
Exact anchored right congruence and finite transfer matrices; conjecture open & E\\
Dyadic numerators &
Autonomous single-integer decoder; $(N_n)$ is not $P$-recursive & L/E\\
Mellin boundary &
$\Re s=0$ is natural; all Gumbel moments are recovered from $Z_N$ & A/E\\
Algebraic separation &
$F_N(r)$ is algebraically independent from $\pi$; relative $\sqrt N$ jet bound & L/E\\
Alternating prime spectrum &
Critical rank-one model, exact right edge, and vague Jessen zero-measure limit & A/L/E\\
Global parity &
Positive lower densities, thick excursions, and punctured Fourier convergence; $o(x)$ open & L/E\\
Balanced determinant &
Exact trace field of transcendence degree $2$ and rational-fibre formula & L/E\\
Prime-parity word algebra &
Locally nilpotent nonnilpotent radical with $\dim\mathcal J_n=p_b(n)-2$ & E\\
\bottomrule
\end{tabularx}
\end{center}

\appendix

\section{Quantitative truncated automaticity}

Fix a base $q\ge2$ and use the least-significant-digit-first $q$-adic convention in
Dubbe's definition \cite{Dubbe}.  Let $A_{\PP}^{(q)}(x)$ be the least number of states in
an automaton recognizing primality correctly for all integers up to $x$, and let
$A_b^{(q)}(x)$ be the analogous minimum for a DFAO producing $b_n$ on that range.

\begin{proposition}[Transfer of Dubbe's lower bound]\label{prop:Dubbe-transfer}
\status{L+E}
For every fixed $q$ there is $C_q>0$ such that
\begin{equation}\label{eq:Dubbe-transfer}
 A_b^{(q)}(x)
 \gg_q x^{1/2}
 \exp\!\left[-C_q(\log\log x)^2\log\log\log x\right]
\end{equation}
for all sufficiently large $x$.
\end{proposition}

\begin{proof}
Suppose a DFAO with $S$ states produces $b_n$ through the required range.  Run one copy
on $n$ and a second copy behind a finite subtract-one/borrow transducer on $n-1$; take
the XOR of the two outputs.  By \eqref{eq:xor} the product machine recognizes primes and
has $O_q(S^2)$ states.  Therefore
\[
 A_{\PP}^{(q)}(x)\ll_q\bigl(A_b^{(q)}(x+1)\bigr)^2.
\]
Dubbe's lower bound
\[
 A_{\PP}^{(q)}(x)
 \ge x\exp\!\left[-c_q(\log\log x)^2\log\log\log x\right]
\]
gives \eqref{eq:Dubbe-transfer} after taking square roots and changing the constant.
\end{proof}

The $O_q(S^2)$ formulation is representation-safe.  An exact $2S^2+O(1)$ count
requires a fixed high-order zero-padding convention so that subtracting one cannot create
an unhandled terminal leading zero.

\section{Notation and reproducibility checklist}

\begin{center}
\small
\begin{tabularx}{\textwidth}{@{}l>{\raggedright\arraybackslash}X@{}}
\toprule
Symbol & Meaning\\
\midrule
$p_k,g_k$ & $k$-th prime and the consecutive gap $p_{k+1}-p_k$\\
$\chi_n,b_n,\sigma_n$ & prime indicator, cumulative prime parity, and its sign\\
$A(z),B(z)$ & signed-prime difference series and cumulative-parity series\\
$C_{PP}$ & $\sum_k(-1)^{k+1}2^{-p_k}$; its digit word belongs to $2C_{PP}$\\
$C_X(h)$ & prime-pair correlation weighted by prime-index difference parity\\
$D_X(h),\mathcal F_{X,h}(z)$ & parity-overlap correlation and prime-rank phase polynomial\\
$\widehat\mu_N,\widehat\nu_N$ & finite diffraction and its renormalized high-pass measure\\
$x_k,\delta_k$ & prime-time doubling point and its nearest-endpoint distance\\
$z_k$ & corrected Gauss tail $G^k(\alpha_P)=[0;g_k,g_{k+1},\ldots]$\\
$q_n,Q_m$ & full-gap convergent denominator and $Q_m=q_{2m}$\\
$H_n$ & half-gap continuant denominator\\
$N_n$ & odd numerator of $C_n=N_n/2^{p_n}$\\
$F_Q,F_H,F_N$ & support series with exponents $Q_m,H_n,N_n$\\
$Z_N,A_m$ & numerator Dirichlet series and its standard-Gumbel boundary moments\\
$\Palt,\Psi$ & alternating prime-index Dirichlet series and its perturbation determinant\\
$\mathsf A_{\rm o},\mathsf A_{\rm e}$ & compact operators carrying the odd- and even-index inverse-prime spectra\\
$J,J''$ & Jessen function and its distributional zero-counting measure\\
\bottomrule
\end{tabularx}
\end{center}

All exact finite identities can be checked from a finite prime list:
\begin{enumerate}
 \item build $b_n$ and verify \eqref{eq:xor}, \eqref{eq:twin-curvature}, and
 \eqref{eq:cumulative-anti};
 \item iterate \eqref{eq:q-recurrence} and compare the three floors in
 \eqref{eq:three-decoders};
 \item iterate the half-gap recurrence and verify \eqref{eq:H-mod3};
 \item iterate \eqref{eq:N-recurrence} and verify the valuation
 \eqref{eq:2adic-decoder};
 \item implement the anchored Gilbreath state using \eqref{eq:boundary-update} and
 \eqref{eq:anchored-residue-update}, and verify the counterexample in
 \cref{rem:bisimulation-counterexample}.
\end{enumerate}
Such computations are useful for detecting finite-identity errors, but they do not replace
the limiting arguments or the cited analytic theorems.

\section{A balanced Gamma--Liouville determinant}

Fix $r\in\mathbb Q$ with $0<|r|<1$ and set $L=L_r$.  This appendix records a
deliberately constructed bridge between this Liouville value and the Gamma trace tower.
It is a consequence of the preceding independence theorem, not a new general principle
in operator theory.  Put
\[
 \mathsf C_r=[L]\oplus\mathsf A_\Gamma\oplus(-\mathsf A_\Gamma).
\]
Its nonexceptional eigenvalues are $1/n$ and $-1/n$, each with multiplicity five.

\begin{theorem}[Balanced trace tower and determinant]
\label{thm:balanced-Gamma-Liouville}
\status{E+L}
$\mathsf C_r$ is compact, self-adjoint and Hilbert--Schmidt, but not trace class.  Its
balanced first trace and ordinary higher traces are
\begin{equation}\label{eq:balanced-traces}
\begin{aligned}
 T_1:=\Tr_{\rm pv}\mathsf C_r
 &:=\lim_{N\to\infty}
 \Tr\bigl([L]\oplus\mathsf A_{\Gamma,N}\oplus(-\mathsf A_{\Gamma,N})\bigr)=L,\\
 T_m:=\Tr\mathsf C_r^m
 &=\begin{cases}
  L^m,&m\text{ odd},\\
  L^m+10\zeta(m),&m\text{ even}
 \end{cases}
 \qquad(m\ge2).
\end{aligned}
\end{equation}
Every $T_m$ is transcendental,
\[
 \overline{\mathbb Q}(T_1,T_2,\ldots)
 =\overline{\mathbb Q}(L,\pi^2),
 \qquad \operatorname{trdeg}_{\overline{\mathbb Q}}=2,
\]
and two distinct traces are algebraically dependent exactly when both indices are odd.

The balanced regularized determinant is
\begin{equation}\label{eq:balanced-determinant}
 \Delta_L(t):=e^{-tL}\det\nolimits_2(I-t\mathsf C_r)
 =(1-tL)\left(\frac{\sin\pi t}{\pi t}\right)^5.
\end{equation}
Every nonconstant Taylor coefficient is transcendental.  Its zero divisor consists of
the nonzero integers, each with multiplicity five, and the simple Liouville zero $1/L$.
If distinct $q_1,\ldots,q_s\in\mathbb Q\setminus\mathbb Z$, then
\begin{equation}\label{eq:balanced-rational-fibres}
 \operatorname{trdeg}_{\overline{\mathbb Q}}
 \overline{\mathbb Q}(\Delta_L(q_1),\ldots,\Delta_L(q_s))=\min\{s,2\}.
\end{equation}
\end{theorem}

\begin{proof}
The spectral symmetry gives \eqref{eq:balanced-traces}.  Since $L$ and $\pi$ are
algebraically independent by \cref{thm:FN-pi-independence}, all transcendence and field
claims follow; for two even traces, the Jacobian of
$x^a+c_ay^a$ and $x^b+c_by^b$ is nonzero when $a\ne b$.

In the determinant product, exponential factors cancel between each $\pm1/n$ pair and
the exceptional eigenvalue contributes $(1-tL)e^{tL}$.  Euler's sine product proves
\eqref{eq:balanced-determinant}.  If
$(\sin x/x)^5=\sum_{k\ge0}q_kx^{2k}$, then
$q_k\in\mathbb Q^\times$ (equivalently, these are nonzero even moments of a sum of five
uniform variables, with alternating sign).  Hence the even coefficients are
$q_k\pi^{2k}$ and the odd ones are $-q_kL\pi^{2k}$.

For rational $q\notin\mathbb Z$, set
$a_q=(\sin(\pi q)/q)^5\in\overline{\mathbb Q}^{\times}$ and
$x_q=\Delta_L(q)/a_q=\pi^{-5}(1-qL)$.  For distinct $q_1,q_2$,
\[
 \pi^{-5}=\frac{q_2x_{q_1}-q_1x_{q_2}}{q_2-q_1},
 \qquad
 L=\frac{x_{q_1}-x_{q_2}}{q_2x_{q_1}-q_1x_{q_2}}.
\]
Thus any pair recovers the independent generators and proves
\eqref{eq:balanced-rational-fibres}; one value is already transcendental.
\end{proof}

\section{The radical of the prime-parity word algebra}

Let $K$ be a field and let $\operatorname{Fac}(b)$ be the factor language of the
prime-parity word.  Form the monomial algebra
\begin{equation}\label{eq:word-algebra}
 \mathcal A_b:=K\langle x_0,x_1\rangle/
 \langle w:w\notin\operatorname{Fac}(b)\rangle,
\end{equation}
and let $\mathcal J$ be the ideal spanned by surviving monomials containing both letters.

\begin{theorem}[Locally nilpotent complexity radical]
\label{thm:word-algebra-radical}
\status{E}
\begin{equation}\label{eq:word-algebra-radical}
 \operatorname{Jac}(\mathcal A_b)=\mathcal J,
 \qquad
 \mathcal A_b/\mathcal J\cong K[x,y]/(xy).
\end{equation}
The ideal $\mathcal J$ is locally nilpotent but not nilpotent, and its homogeneous
pieces satisfy
\begin{equation}\label{eq:radical-dimension}
 \dim_K\mathcal J_n=p_b(n)-2.
\end{equation}
Thus all but two dimensions of the superpolynomial factor growth lie inside the
Jacobson radical.
\end{theorem}

\begin{proof}
Each mixed surviving word contains a transition.  Given finitely many generators of
$\mathcal J$, any surviving product of $m$ of them has at least $m$ internal
transitions and length at most $Cm$ for a fixed $C$.  This contradicts the upper Banach
density zero of the prime transition set when $m$ is large.  Hence $\mathcal J$ is
locally nilpotent and therefore lies in the Jacobson radical.  It is not nilpotent,
because arbitrarily long finite portions of $b$ contain arbitrarily many prime
transitions and can be split into arbitrarily many mixed factors.

Modulo $\mathcal J$, only pure powers remain, giving $K[x,y]/(xy)$, a semiprimitive
algebra.  Hence the radical is exactly $\mathcal J$.  It remains to justify the two pure
words in every degree.  Choose distinct primes $q_j>|j|$ for
$0<|j|\le n$ and solve $P\equiv-j\pmod{q_j}$.  The resulting residue is coprime to
$Q=\prod q_j$, so Dirichlet's theorem supplies a prime $P$ in that progression; every
$P+j$ with $0<|j|\le n$ is composite.  The long constant blocks on the two sides of
$P$ have opposite bits, so both $0^n$ and $1^n$ occur.  Exactly these two degree-$n$
basis monomials lie outside $\mathcal J$, proving \eqref{eq:radical-dimension}.
For the general word-algebra framework, compare \cite{BellWordAlgebra}.
\end{proof}

\begin{thebibliography}{99}

\bibitem{Alboiu}
M.~Alboiu,
\emph{On the parity of the prime-counting function and related problems},
Ramanujan J. \textbf{38} (2015), 179--187.
\href{https://doi.org/10.1007/s11139-014-9596-1}{doi:10.1007/s11139-014-9596-1}.

\bibitem{AlloucheShallit}
J.-P.~Allouche and J.~Shallit,
\emph{Automatic Sequences: Theory, Applications, Generalizations},
Cambridge University Press, 2003.

\bibitem{BayartKouroupis}
F.~Bayart and A.~Kouroupis,
\emph{On universality of general Dirichlet series},
Constr. Approx. (2026).
\href{https://doi.org/10.1007/s00365-026-09772-5}{doi:10.1007/s00365-026-09772-5};
\href{https://arxiv.org/abs/2402.00656}{arXiv:2402.00656}.

\bibitem{Bugeaud}
Y.~Bugeaud,
\emph{Approximation by Algebraic Numbers},
Cambridge Tracts in Mathematics 160, Cambridge University Press, 2004.

\bibitem{Carlson}
F.~Carlson,
\emph{\"Uber Potenzreihen mit ganzzahligen Koeffizienten},
Math. Z. \textbf{9} (1921), 1--13.
\href{https://doi.org/10.1007/BF01378331}{doi:10.1007/BF01378331}.

\bibitem{CHT}
Z.~Chase, Z.~Hunter and T.~Tao,
\emph{Gilbreath's conjecture: a Cram\'er random model and a deterministic analysis},
\href{https://arxiv.org/abs/2607.08712}{arXiv:2607.08712} (2026).

\bibitem{ChristolEtAl}
G.~Christol, T.~Kamae, M.~Mend\`es France and G.~Rauzy,
\emph{Suites alg\'ebriques, automates et substitutions},
Bull. Soc. Math. France \textbf{108} (1980), 401--419.
\href{https://doi.org/10.24033/bsmf.1926}{doi:10.24033/bsmf.1926}.

\bibitem{ChungLi}
P.~N.~Chung and S.~Li,
\emph{On the residue classes of $\pi(n)$ modulo $t$},
Integers \textbf{13} (2013), Paper A79.
\href{https://doi.org/10.5281/zenodo.10409040}{doi:10.5281/zenodo.10409040}.

\bibitem{CorvajaZannier}
P.~Corvaja and U.~Zannier,
\emph{Some new applications of the Subspace Theorem},
Compos. Math. \textbf{131} (2002), 319--340.
\href{https://doi.org/10.1023/A:1015594913393}{doi:10.1023/A:1015594913393}.

\bibitem{Devyatov}
R.~A.~Devyatov,
\emph{On subword complexity of morphic sequences},
in \emph{Computer Science---Theory and Applications}, LNCS 5010,
Springer, 2008, 146--157; see also
\href{https://arxiv.org/abs/1502.02310}{arXiv:1502.02310}.

\bibitem{Dubbe}
T.~Dubbe,
\emph{The automaticity of the set of primes},
\href{https://arxiv.org/abs/2409.04314}{arXiv:2409.04314} (2024).

\bibitem{FGKMT}
K.~Ford, B.~Green, S.~Konyagin, J.~Maynard and T.~Tao,
\emph{Long gaps between primes},
J. Amer. Math. Soc. \textbf{31} (2018), 65--105.
\href{https://doi.org/10.1090/jams/876}{doi:10.1090/jams/876}.

\bibitem{HartmanisShank}
J.~Hartmanis and H.~Shank,
\emph{On the recognition of primes by automata},
J. ACM \textbf{15} (1968), 382--389.
\href{https://doi.org/10.1145/321466.321470}{doi:10.1145/321466.321470}.

\bibitem{HWY}
W.~Huang, Z.~Wang and X.~Ye,
\emph{Measure complexity and M\"obius disjointness},
Adv. Math. \textbf{347} (2019), 827--858.
\href{https://doi.org/10.1016/j.aim.2019.03.007}{doi:10.1016/j.aim.2019.03.007}.

\bibitem{IK}
H.~Iwaniec and E.~Kowalski,
\emph{Analytic Number Theory},
AMS Colloquium Publications 53, 2004.

\bibitem{MaynardSmallGaps}
J.~Maynard,
\emph{Small gaps between primes},
Ann. of Math. \textbf{181} (2015), 383--413.
\href{https://doi.org/10.4007/annals.2015.181.1.7}{doi:10.4007/annals.2015.181.1.7}.

\bibitem{MontgomeryVaughan}
H.~L.~Montgomery and R.~C.~Vaughan,
\emph{The large sieve},
Mathematika \textbf{20} (1973), 119--134.
\href{https://doi.org/10.1112/S0025579300004708}{doi:10.1112/S0025579300004708}.

\bibitem{MorseHedlund}
M.~Morse and G.~A.~Hedlund,
\emph{Symbolic dynamics},
Amer. J. Math. \textbf{60} (1938), 815--866.

\bibitem{OEIS}
OEIS Foundation Inc.,
\emph{The On-Line Encyclopedia of Integer Sequences},
entries A071986, A109139, A109140, A122150 and A122153,
accessed 1 September 2026.
\url{https://oeis.org}.

\bibitem{Polymath8b}
D.~H.~J.~Polymath,
\emph{Variants of the Selberg sieve, and bounded intervals containing many primes},
Res. Math. Sci. \textbf{1} (2014), Article 12; Erratum \textbf{2} (2015), Article 15.
\href{https://doi.org/10.1186/s40687-014-0012-7}{doi:10.1186/s40687-014-0012-7}.

\bibitem{Polya}
G.~P\'olya,
\emph{\"Uber Potenzreihen mit ganzzahligen Koeffizienten},
Math. Ann. \textbf{77} (1916), 497--513.
\href{https://doi.org/10.1007/BF01456965}{doi:10.1007/BF01456965}.

\bibitem{Selberg}
A.~Selberg,
\emph{The general sieve-method and its place in prime number theory},
Proc. ICM Cambridge, Mass. 1950, Vol.~1, AMS, 1952, 286--292.

\bibitem{SepulcreVidal}
J.~M.~Sepulcre and T.~Vidal,
\emph{On the real projections of zeros of analytic almost periodic functions},
Carpathian J. Math. \textbf{38} (2022), 489--501.
\href{https://doi.org/10.37193/CJM.2022.02.18}{doi:10.37193/CJM.2022.02.18};
\href{https://arxiv.org/abs/1805.02041}{arXiv:1805.02041}.

\bibitem{Stanley}
R.~P.~Stanley,
\emph{Differentiably finite power series},
European J. Combin. \textbf{1} (1980), 175--188.
\href{https://doi.org/10.1016/S0195-6698(80)80051-5}{doi:10.1016/S0195-6698(80)80051-5}.

\bibitem{TaoErdos}
T.~Tao,
\emph{The convergence of an alternating series of Erd\H{o}s, assuming the
Hardy--Littlewood prime tuples conjecture},
Commun. Amer. Math. Soc. \textbf{4} (2024), 80--96.
\href{https://doi.org/10.1090/cams/29}{doi:10.1090/cams/29}.

\bibitem{Wei}
F.~Wei,
\emph{Disjointness of M\"obius from asymptotically periodic functions},
Pure Appl. Math. Q. \textbf{18} (2022), 863--922.
\href{https://doi.org/10.4310/PAMQ.2022.v18.n3.a3}{doi:10.4310/PAMQ.2022.v18.n3.a3}.

\bibitem{AdamczewskiBellSmertnig}
B.~Adamczewski, J.~P.~Bell and D.~Smertnig,
\emph{A height gap theorem for coefficients of Mahler functions},
J. Eur. Math. Soc. \textbf{25} (2023), 2525--2571.
\href{https://doi.org/10.4171/JEMS/1244}{doi:10.4171/JEMS/1244}.

\bibitem{BellWordAlgebra}
J.~P.~Bell,
\emph{Topological invariants for words of linear factor complexity},
\href{https://arxiv.org/abs/2202.00643}{arXiv:2202.00643} (2022).

\bibitem{BirmanYafaevSSF}
M.~Sh.~Birman and D.~R.~Yafaev,
\emph{The spectral shift function: the work of M.~G.~Krein and its further development},
St. Petersburg Math. J. \textbf{4} (1993), 833--870.
\url{https://www.mathnet.ru/eng/aa343}.

\bibitem{BorchseniusJessen1948}
V.~Borchsenius and B.~Jessen,
\emph{Mean motions and values of the Riemann zeta function},
Acta Math. \textbf{80} (1948), 97--166.
\href{https://doi.org/10.1007/BF02393647}{doi:10.1007/BF02393647}.

\bibitem{ChalebgwaMorris}
T.~P.~Chalebgwa and S.~A.~Morris,
\emph{Sin, cos, exp and log of Liouville numbers},
Bull. Aust. Math. Soc. \textbf{108} (2023), 81--85.
\href{https://doi.org/10.1017/S000497272200140X}
{doi:10.1017/S000497272200140X}.

\bibitem{IonascuRankOne}
E.~J.~Ionascu,
\emph{Rank-one perturbations of diagonal operators},
Integral Equations Operator Theory \textbf{39} (2001), 421--440.
\href{https://doi.org/10.1007/BF01203323}{doi:10.1007/BF01203323}.

\bibitem{JessenTornehave1945}
B.~Jessen and H.~Tornehave,
\emph{Mean motions and zeros of almost periodic functions},
Acta Math. \textbf{77} (1945), 137--279.
\href{https://doi.org/10.1007/BF02392225}{doi:10.1007/BF02392225}.

\bibitem{Landau1909}
E.~Landau,
\emph{Handbuch der Lehre von der Verteilung der Primzahlen}, Vol.~II,
B.~G.~Teubner, Leipzig, 1909.

\bibitem{LenzSpindelerStrungaru}
D.~Lenz, T.~Spindeler and N.~Strungaru,
\emph{Pure point diffraction and mean, Besicovitch and Weyl almost periodicity},
\href{https://arxiv.org/abs/2006.10821}{arXiv:2006.10821}, revised 2026.

\bibitem{LiawRankOne}
C.~Liaw,
\emph{Rank one and finite rank perturbations},
\href{https://arxiv.org/abs/1205.4376}{arXiv:1205.4376} (2012).

\bibitem{LipshitzRubel}
L.~Lipshitz and L.~A.~Rubel,
\emph{A gap theorem for power series solutions of algebraic differential equations},
Amer. J. Math. \textbf{108} (1986), 1193--1214.
\href{https://doi.org/10.2307/2374602}{doi:10.2307/2374602}.

\bibitem{MRTChowla}
K.~Matom\"aki, M.~Radziwi\l\l{} and T.~Tao,
\emph{An averaged form of Chowla's conjecture},
Algebra Number Theory \textbf{9} (2015), 2167--2196.
\href{https://doi.org/10.2140/ant.2015.9.2167}{doi:10.2140/ant.2015.9.2167};
corrected version \href{https://arxiv.org/abs/1503.05121v3}{arXiv:1503.05121v3}.

\bibitem{Murai}
T.~Murai,
\emph{The value-distribution of lacunary series and a conjecture of Paley},
Ann. Inst. Fourier \textbf{31} (1981), no.~1, 135--156.
\href{https://doi.org/10.5802/aif.820}{doi:10.5802/aif.820}.

\bibitem{PowersGamma}
M.~R.~Powers,
\emph{``Infinitely often'' transcendence of Gamma-function derivatives},
\href{https://arxiv.org/abs/2601.18474}{arXiv:2601.18474}, version 4 (2026).

\bibitem{SchillingBernstein}
R.~L.~Schilling, R.~Song and Z.~Vondra\v{c}ek,
\emph{Bernstein Functions: Theory and Applications}, 2nd ed.,
De Gruyter Studies in Mathematics 37, De Gruyter, 2012.

\bibitem{SimonTraceIdeals}
B.~Simon,
\emph{Trace Ideals and Their Applications}, 2nd ed.,
Mathematical Surveys and Monographs 120, American Mathematical Society, 2005.
\href{https://doi.org/10.1090/surv/120}{doi:10.1090/surv/120}.

\end{thebibliography}

\end{document}
